\documentclass[11pt, a4paper]{article}

% --- Geometría y Cabeceras ---
\usepackage[margin=1in, headheight=14pt]{geometry}
\usepackage{fancyhdr}
\pagestyle{fancy}
\fancyhf{}
\fancyhead[L]{\textit{Teoría de Grafos}}
\fancyhead[R]{\textit{Bustos Jordi}}
\fancyfoot[C]{\thepage}
\renewcommand{\headrulewidth}{0.4pt}

% --- Matemática e Idioma ---
\usepackage{amsmath,amsthm,amssymb}
\usepackage[spanish, es-tabla]{babel}

% --- Tipografía para XeLaTeX ---
% Libertinus es excelente para matemática moderna y muy estética.
\usepackage{unicode-math}

% --- Cajas y Colores ---
\usepackage[dvipsnames]{xcolor}
\usepackage[many]{tcolorbox}
\usepackage{enumitem}
\usepackage{graphicx}

% --- TikZ y Estilos para Grafos ---
\usepackage{tikz}
\usetikzlibrary{calc, positioning, arrows.meta}

% Definimos estilos globales para que todos los grafos se vean iguales
\tikzset{
    % Estilo para vértices normales
    vertex/.style={circle, draw=black!80, fill=gray!10, thick, inner sep=2pt, minimum size=6mm, font=\small\sffamily},
    % Estilo para aristas
    edge/.style={thick, draw=black!70},
    % Estilo para grafos dirigidos
    directed/.style={edge, ->, >={Stealth[length=2.5mm]}},
    % Vértices coloreados para 2-coloreos
    vertex A/.style={vertex, fill=blue!15, draw=blue!70!black},
    vertex B/.style={vertex, fill=red!15, draw=red!70!black}
}

% --- Algoritmos ---
\usepackage{algorithm}
\usepackage{algpseudocode}
\usepackage{xparse}
\usepackage{etoolbox}

% --- Entornos Matemáticos ---
\theoremstyle{plain}
\newtheorem{proposicion}{Proposición}[section]
\newtheorem{teorema}{Teorema}[section]

\theoremstyle{definition}
\newtheorem{definition}{Definición}[section]

% --- Entorno de Ejercicio Estético ---
\newtcolorbox{ejercicio}[1]{
    colback=blue!2!white,      % Fondo súper sutil
    colframe=blue!40!black,    % Borde elegante
    fonttitle=\bfseries,
    title=Ejercicio #1,
    arc=3mm,                   % Bordes redondeados
    boxrule=0.8pt,
    left=4pt, right=4pt, top=4pt, bottom=4pt
}

% --- Macros y Atajos ---
\newcommand{\R}{\mathbb{R}}
\newcommand{\Z}{\mathbb{Z}}
\newcommand{\N}{\mathbb{N}}
\newcommand{\Q}{\mathbb{Q}}
\newcommand{\C}{\mathbb{C}}

\setlist[enumerate,1]{label=(\alph*), leftmargin=*, itemsep=2pt}

% Comando de colores para etiquetas en texto
\newcommand{\vlab}[1]{%
    \ifstrequal{#1}{A}{\textcolor{blue!70!black}{\textbf{A}}}{%
    \ifstrequal{#1}{B}{\textcolor{red!70!black}{\textbf{B}}}{#1}}%
}

% Un paso del algoritmo BFS del Ejercicio 11 sobre el grafo de Petersen: recibe
% las etiquetas de los 5 vértices exteriores o1,...,o5 (#1) y de los 5
% vértices interiores i1,...,i5 (#2), como listas separadas por comas, y
% dibuja el grafo con esas etiquetas, evitando repetir el mismo dibujo cuatro veces.
\NewDocumentCommand{\PetersenPaso}{ m m }{%
    \begin{tikzpicture}[vertex/.style={circle, draw, fill=black, inner sep=1.5pt, outer sep=0pt}]
        \def\R{2.6}
        \def\r{1.2}

        \foreach [count=\i] \lab in {#1} {
            \ifcase\i\relax
            \or \node[vertex] (o1) at (90:\R) {}; \node[above=3pt] at (o1) {\vlab{\lab}};
            \or \node[vertex] (o2) at (18:\R) {}; \node[above right=1pt] at (o2) {\vlab{\lab}};
            \or \node[vertex] (o3) at (306:\R) {}; \node[below right=1pt] at (o3) {\vlab{\lab}};
            \or \node[vertex] (o4) at (234:\R) {}; \node[below left=1pt] at (o4) {\vlab{\lab}};
            \or \node[vertex] (o5) at (162:\R) {}; \node[above left=1pt] at (o5) {\vlab{\lab}};
            \fi
        }
        \foreach [count=\i] \lab in {#2} {
            \ifcase\i\relax
            \or \node[vertex] (i1) at (90:\r) {}; \node[right=2pt] at (i1) {\vlab{\lab}};
            \or \node[vertex] (i2) at (18:\r) {}; \node[below right=1pt] at (i2) {\vlab{\lab}};
            \or \node[vertex] (i3) at (306:\r) {}; \node[left=2pt] at (i3) {\vlab{\lab}};
            \or \node[vertex] (i4) at (234:\r) {}; \node[below right=1pt] at (i4) {\vlab{\lab}};
            \or \node[vertex] (i5) at (162:\r) {}; \node[above right=1pt] at (i5) {\vlab{\lab}};
            \fi
        }

        \draw[thick] (o1) -- (o2) -- (o3) -- (o4) -- (o5) -- (o1);
        \draw[thick] (i1) -- (i3) -- (i5) -- (i2) -- (i4) -- (i1);
        \foreach \x in {1,...,5} {
                \draw[thick] (o\x) -- (i\x);
            }
    \end{tikzpicture}%
}

\newtcolorbox{ejercicioEntregar}[1]{
    colback=teal!5!white,       % Fondo verde/azulado muy sutil
    colframe=teal!60!black,     % Borde más oscuro y distintivo
    fonttitle=\bfseries,
    % El \hfill empuja el texto "Para entregar" hacia la derecha de la caja
    title={Ejercicio #1 \hfill \normalfont\itshape \(\star\) Seleccionado para entregar},
    arc=3mm,
    boxrule=1.2pt,              % Borde un pelín más grueso para que destaque
    left=4pt, right=4pt, top=4pt, bottom=4pt
}

% --- Hipervínculos ---
\usepackage[colorlinks=true, linkcolor=blue!60!black, urlcolor=blue!60!black, citecolor=blue!60!black]{hyperref}

\begin{document}

% --- Portada Personalizada ---
\begin{titlepage}
	\centering
	\vspace*{4cm}
	{\Huge\bfseries Teoría de Grafos\par}
	\vspace{1.5cm}
	{\Large Resolución de Prácticas I -- VI\par}
	\vspace{1.5cm}
	{\large \textbf{Autor:}\\ Bustos Jordi\par}
	\vfill
	{\large Universidad Nacional de La Plata\par}
	{\large \today\par}
\end{titlepage}

\clearpage
\tableofcontents
\clearpage

\section{Práctica I}
\begin{definition}
    Un grafo \( G \) es simple \( \iff \) cada par de vértices distintos está conectado por como máximo una arista y no tiene lazos.
\end{definition}

\begin{ejercicio}{1}
    Probar que, si \( G \) es un grafo simple, el número de aristas de \( G \) es menor o igual que \( n(n-1)/2 \), siendo \( n \) el número de vértices.
\end{ejercicio}
\begin{proof}
    Al ser \( G \) un grafo simple, no puede tener aristas múltiples ni lazos. Por lo tanto, cada par de vértices distintos puede estar conectado por
    una arista como máximo. El número de pares distintos de vértices en un conjunto de \( n \) vértices está dado por la combinación \[
        \binom{n}{2} = \frac{n(n-1)}{2}.
    \]
    que representa el número máximo de aristas posibles en un grafo simple con \( n \) vértices y, por lo tanto, el número de aristas de \( G \) es menor o igual que \( n(n-1)/2 \).
\end{proof}

\vspace{0.5cm}

\begin{definition}\label{def:grado}
    Definimos el grado de un vértice \( v \) en un grafo \( G \) como el número de aristas incidentes a \( v \). Se denota por \( \deg(v) \).
\end{definition}

\begin{ejercicio}{2}
    Dado \( n \) entero mayor o igual que 1, hallar una fórmula para la cantidad de grafos simples con conjunto de vértices \( \{1,\dots, n\} \). ¿En cuántos de esos grafos simples el grado del vértice 1 es menor o igual que 2? (la respuesta debe ser una expresión en función de \( n \) y completamente simplificada)
\end{ejercicio}
\begin{proof}
    Sea \( n \) el número de vértices. Un grafo simple con \( n \) vértices puede tener o no aristas entre cada par de vértices distintos. Dado que hay \( \binom{n}{2} \) números de pares distintos de vértices, y cada par puede estar conectado por una arista o no, el número total de grafos simples con \( n \) vértices es de \( 2^{\binom{n}{2}} \).

    Para la segunda parte, consideremos el vértice 1. El grado del vértice 1 es el número de aristas que lo conectan con los otros \( n-1 \) vértices. Para que su grado sea menor o igual que 2, podemos tener las siguientes posibilidades:
    \begin{enumerate}
        \item Grado 0: Ninguna arista conectando el vértice 1 con los otros \( n-1 \) vértices. Esto da \( 2^{\binom{n-1}{2}} \) formas.
        \item Grado 1: Una arista conectando el vértice 1 con uno de los \( n-1 \) vértices. Esto da \( (n-1) \cdot 2^{\binom{n-1}{2}} \) formas.
        \item Grado 2: Dos aristas conectando el vértice 1 con dos de los \( n-1 \) vértices. Esto da \( \binom{n-1}{2} \cdot 2^{\binom{n-1}{2}} \) formas.
    \end{enumerate}
    Luego, el número total de grafos simples con el grado del vértice 1 menor o igual que 2 es:
    \[
        \left [ 1 + (n-1) + \binom{n-1}{2} \right] \cdot 2^{\binom{n-1}{2}}
    \]
\end{proof}

\begin{definition}
    Un grafo \( G \) se dice \textbf{conexo} si, para todo par de vértices \( u \) y \( v \) de \( G \), existe una sucesión de vértices que comienza en \( u \), termina en \( v \) y tal que vértices consecutivos son adyacentes. Caso contrario, diremos que \( G \) es \textbf{disconexo}.
\end{definition}

\begin{definition}\label{def:componente-conexa}
    Una \textbf{componente conexa} de un grafo \( G \) es un subgrafo conexo maximal de \( G \), es decir, un subgrafo conexo que no está contenido en ningún subgrafo conexo más grande de \( G \).
\end{definition}

\begin{ejercicio}{3}
    Analizar si es verdadero o falso: el complemento de un grafo simple disconexo es conexo.
\end{ejercicio}
\begin{proof}
    \emph{Claim}: El complemento de un grafo simple disconexo es conexo.

    Sea \( G \) un grafo simple disconexo y sea \( \overline{G} \) su complemento. Como \( G \) es disconexo, sus vértices se reparten en al menos dos componentes conexas distintas. Veamos que \( \overline{G} \) es conexo, es decir, que para cualquier par de vértices \( x \) e \( y \) de \( G \) existe una sucesión de vértices adyacentes en \( \overline{G} \) que los conecta. Puede ocurrir:
    \begin{enumerate}
        \item \( x \) e \( y \) pertenecen a componentes distintas de \( G \). Como no puede haber aristas entre vértices de componentes distintas de \( G \), \( x \) e \( y \) no son adyacentes en \( G \). Luego la arista \( xy \) sí pertenece a \( \overline{G} \), y \( x \) e \( y \) quedan conectados directamente en \( \overline{G} \).
        \item \( x \) e \( y \) pertenecen a la misma componente de \( G \). Como \( G \) es disconexo, existe otra componente distinta de la cual tomamos un vértice \( z \). Por el caso anterior aplicado a los pares \( x,z \) y \( z,y \), las aristas \( xz \) y \( zy \) pertenecen a \( \overline{G} \). Así, \( x - z - y \) es una sucesión de vértices adyacentes en \( \overline{G} \) que conecta a \( x \) con \( y \).
    \end{enumerate}
    En ambos casos hallamos una sucesión de vértices adyacentes en \( \overline{G} \) que conecta a \( x \) con \( y \). Como \( x \) e \( y \) eran arbitrarios, concluimos que \( \overline{G} \) es conexo. Por lo tanto, la afirmación es verdadera.
\end{proof}

\vspace{0.5cm}

\begin{definition}
    Un \textbf{clique} en un grafo \( G \) es un conjunto de vértices de \( G \) tal que cada par de vértices distintos en el conjunto está conectado por una arista en \( G \). En otras palabras, un clique es un subgrafo completo de \( G \).
\end{definition}

\begin{definition}\label{def:conjunto-independiente}
    Un \textbf{conjunto independiente} en un grafo \( G \) es un conjunto de vértices de \( G \) tal que ningún par de vértices distintos en el conjunto está conectado por una arista en \( G \). En otras palabras, un conjunto independiente es un subgrafo sin aristas.
\end{definition}

\begin{ejercicio}{4}
    Determinar el tamaño máximo de un clique y el tamaño máximo de un conjunto independiente para el grafo de abajo. Justificar por qué no puede haber un clique o un conjunto independiente de tamaño mayor.
    \begin{center}
        \begin{tikzpicture}[every node/.style={circle, draw, fill=black, inner sep=1.5pt}]
            \node (n) at (0, 1) {};
            \node (s) at (0, -1) {};
            \node (c1) at (-1, 0) {};
            \node (c2) at (-0.5, 0) {};
            \node (c3) at (0.5, 0) {};
            \node (c4) at (1, 0) {};

            \draw[thick] (n) -- (c1) -- (s) -- (c4) -- (n);
            \draw[thick] (n) -- (c2) -- (s) -- (c3) -- (n);
            \draw[thick] (c1) -- (c2) -- (c3) -- (c4);
            \draw[thick] (n) to[out=-20, in=20, distance=2.5cm] (s);
        \end{tikzpicture}
    \end{center}
\end{ejercicio}
\begin{proof}
    Es claro que para formar un clique maximal necesitamos tomar el polo norte y el polo sur, ya que ambos están conectados con todos los vértices y entre sí. Luego, podemos tomar un vértice del conjunto de 4 vértices que están en el medio, ya que todos ellos están conectados con el polo norte y el polo sur, finalmente podemos tomar un vértice más del conjunto de 4 vértices que están en el medio que sea adyacente al vértice que ya tomamos. Por lo tanto, el tamaño máximo de un clique es 4, pues si agregamos un vértice más del conjunto de 4 vértices que están en el medio, no estará conectado con al menos uno de los vértices que ya tomamos.

    Contrariamente, para formar un conjunto independiente maximal, no podemos tomar ni el polo norte ni el polo sur, ya que ambos están conectados con todos los vértices. Por lo tanto, si elegimos un vértice del conjunto de 4 vértices que están en el medio, no podemos elegir su adyacente, pero sí podemos elegir el otro vértice que está en el medio y que no es adyacente al vértice que ya elegimos. Sin embargo, no podemos elegir ningún otro vértice del centro pues estará conectado a uno de los vértices que ya elegimos. Por lo tanto, el tamaño máximo de un conjunto independiente es 2.
\end{proof}

\begin{ejercicio}{5}
    Probar que cada conjunto de 6 personas contiene al menos 3 mutuamente familiares o 3 mutuamente no familiares.
\end{ejercicio}
\begin{proof}
    Notemos que podemos modelar la situación como un grafo \( G \) donde cada persona es un vértice y hay una arista entre dos vértices si las personas correspondientes son familiares. Entonces, el problema se reduce a demostrar que en cualquier grafo de 6 vértices, existe un clique de tamaño 3 (3 personas mutuamente familiares) o un conjunto independiente de tamaño 3 (3 personas mutuamente no familiares).

    Consideremos un vértice \( v \) en \( G \). Este vértice tiene 5 vecinos posibles. Luego, por el principio del palomar, al menos 3 de estos vecinos deben ser familiares o al menos 3 deben ser no familiares de \( v \). Supongamos que \( v \) tiene al menos 3 vecinos familiares, digamos \( v_1, v_2, v_3 \). Si alguno de estos vértices está conectado entre sí, entonces tenemos un clique de tamaño 3. Si ninguno de ellos está conectado entre sí, entonces \( v_1, v_2, v_3 \) forman un conjunto independiente de tamaño 3. El caso en que \( v \) tiene al menos 3 vecinos no familiares es análogo.
\end{proof}

\vspace{0.5cm}

\begin{definition}
    Sea \( G \) un grafo con conjunto de vértices \( V(G) = \{ v_1, \ldots, v_n \} \) y conjunto de aristas \( E(G) = \{ e_1, \ldots, e_m \} \). La matriz de adyacencia \( A(G) \) es una matriz \( n \times n \) donde el elemento \( a_{ij} \) es 1 si \( v_i \) y \( v_j \) están conectados por una arista y 0 en caso contrario.
\end{definition}

\begin{definition}
    Sea \( G \) un grafo con conjunto de vértices \( V(G) = \{ v_1, \ldots, v_n \} \) y conjunto de aristas \( E(G) = \{ e_1, \ldots, e_m \} \). La matriz de incidencia \( B(G) \) es una matriz \( n \times m \) donde el elemento \( b_{ij} \) es 1 si el vértice \( v_i \) está incidente a la arista \( e_j \) y 0 en caso contrario.
\end{definition}

\begin{ejercicio}{6}
    Sea \( G \) un grafo. Para \( v \in V(G) \) y \( e \in E(G) \), describir las matrices de adyacencia e incidencia de \( G-v \) y \( G-e \) en términos de las correspondientes matrices de \( G \).
\end{ejercicio}
\begin{proof}
    Sea \( V(G) = \{ v_1, \ldots, v_n \} \) y \( E(G) = \{ e_1, \ldots, e_m \} \).

    Luego, si suponemos que \( v = v_i \) para \( i \in \{ 1, \ldots, n \} \) y \( e = e_j = v_p v_q \) para \( j \in \{ 1, \ldots, m \} \), entonces la matriz de adyacencia de \( G - v \) es la matriz \( A(G) \) con la \( i \)-ésima fila y la \( i \)-ésima columna eliminadas. Por otro lado, la matriz de adyacencia de \( G - e \) es la matriz \( A(G) \) con las entradas \( a_{pq} \) y \( a_{qp} \) (correspondientes a los dos vértices incidentes a \( e_j \)) puestas en 0, dejando el resto de la matriz sin cambios, ya que las demás aristas incidentes a \( v_p \) y \( v_q \) siguen presentes en \( G-e \).

    De manera análoga, sea \( B(G) \) la matriz de incidencia de \( G \implies \) la matriz de incidencia de \( G - v \) es la matriz \( B(G) \) con la fila correspondiente a \( v \) eliminada, junto con las columnas correspondientes a las aristas incidentes a \( v \) (pues dichas aristas dejan de existir en \( G - v \)), y la matriz de incidencia de \( G - e \) es la matriz \( B(G) \) con la columna correspondiente a \( e \) eliminada.
\end{proof}

\vspace{0.5cm}

\begin{ejercicio}{7}
    Escribir todas las posibles matrices de adyacencia y de incidencia de \( P_3 \).
\end{ejercicio}
\begin{proof}
    Recordemos que la estructura del grafo camino \( P_3 \) consiste en 3 vértices y 2 aristas, donde siempre hay un vértice central (de grado 2) y dos extremos (de grado 1).
    \vspace{0.5cm}
    \begin{center}
        \begin{tikzpicture}[
                every node/.style={circle, draw, fill=black, inner sep=1.5pt},
            ]
            \node (v1) at (0,0) {};
            \node (v2) at (2,0) {};
            \node (v3) at (4,0) {};
            \draw[thick] (v1) -- (v2) -- (v3);
        \end{tikzpicture}
    \end{center}
    \vspace{0.5cm}
    Las diferentes matrices surgen al permutar las etiquetas de los vértices \( \{v_1, v_2, v_3\} \) y de las aristas \( \{e_1, e_2\} \).

    \textbf{Matrices de adyacencia:} \par
    Dependen únicamente de qué vértice se elija como el vértice central (el de grado 2). Existen \(\binom{3}{1} = 3\) posibilidades:
    \begin{itemize}
        \item Si el centro es \( v_2 \) (el camino es \( v_1 - v_2 - v_3 \)):
              \[ A_1 = \begin{pmatrix} 0 & 1 & 0 \\ 1 & 0 & 1 \\ 0 & 1 & 0 \end{pmatrix} \]
        \item Si el centro es \( v_1 \) (el camino es \( v_2 - v_1 - v_3 \)):
              \[ A_2 = \begin{pmatrix} 0 & 1 & 1 \\ 1 & 0 & 0 \\ 1 & 0 & 0 \end{pmatrix} \]
        \item Si el centro es \( v_3 \) (el camino es \( v_1 - v_3 - v_2 \)):
              \[ A_3 = \begin{pmatrix} 0 & 0 & 1 \\ 0 & 0 & 1 \\ 1 & 1 & 0 \end{pmatrix} \]
    \end{itemize}

    \textbf{Matrices de incidencia:} \par
    Para cada una de las 3 elecciones del vértice central, existen \( 2! = 2 \) formas de ordenar las etiquetas de las aristas en las columnas. Esto nos da un total de 6 matrices posibles:
    \begin{itemize}
        \item Con \( v_2 \) como centro (\( e_1 \) y \( e_2 \) incidentes a \( v_2 \)):
              \[ B_1 = \begin{pmatrix} 1 & 0 \\ 1 & 1 \\ 0 & 1 \end{pmatrix}, \quad B_2 = \begin{pmatrix} 0 & 1 \\ 1 & 1 \\ 1 & 0 \end{pmatrix} \]
        \item Con \( v_1 \) como centro (\( e_1 \) y \( e_2 \) incidentes a \( v_1 \)):
              \[ B_3 = \begin{pmatrix} 1 & 1 \\ 1 & 0 \\ 0 & 1 \end{pmatrix}, \quad B_4 = \begin{pmatrix} 1 & 1 \\ 0 & 1 \\ 1 & 0 \end{pmatrix} \]
        \item Con \( v_3 \) como centro (\( e_1 \) y \( e_2 \) incidentes a \( v_3 \)):
              \[ B_5 = \begin{pmatrix} 1 & 0 \\ 0 & 1 \\ 1 & 1 \end{pmatrix}, \quad B_6 = \begin{pmatrix} 0 & 1 \\ 1 & 0 \\ 1 & 1 \end{pmatrix} \]
    \end{itemize}
\end{proof}

\begin{definition}\label{def:isomorfos}
    Diremos que dos grafos simples \( G \) y \( H \) son isomorfos \( \iff \) existe una biyección \( f : V(G) \to V(H) \) tal que \( \{u,v\} \in E(G) \iff \{f(u), f(v)\} \in E(H)\) para todo \( u,v \in V(G) \).
\end{definition}

\begin{definition}\label{def:complemento}
    Sea \( G \) un grafo, notaremos \( \overline{G} \) al grafo complemento de \( G \), es decir, el grafo con los mismos vértices que \( G \) y con aristas \( \{u,v\} \) tales que \( \{u,v\} \notin E(G) \) para todo \( u,v \in V(G) \).
\end{definition}

\begin{ejercicio}{8}
    Probar que dos grafos simples \( G \) y \( H \) son isomorfos si y sólo si \( \overline{G} \) y \( \overline{H} \) lo son.
\end{ejercicio}
\begin{proof}
    Sea \( f \) el isomorfismo entre \( G \) y \( H \). Entonces, para todo \( u,v \in V(G) \),
    \[
        \{u,v\} \in E(G) \iff \{f(u), f(v)\} \in E(H),
    \]
    lo que implica que \( f \) también es un isomorfismo entre \( \overline{G} \) y \( \overline{H} \).
    Pues, si \( \{ u, v \} \notin E(G) \), debe ser que \( \{ f(u), f(v) \} \notin E(H) \), pues si no \( f \) no cumpliría la condición de isomorfismo. Notemos que \( \{ u, v \} \notin E(G) \iff \{ u, v \} \in E\left(\overline{G}\right) \) para cualquier grafo \( G \).

    La recíproca es análoga.
\end{proof}

\vspace{0.5cm}

\begin{ejercicio}{9}
    Probar que el complemento de cualquiera de estos grafos es isomorfo al otro.
    \vspace{.5cm}
    \begin{center}
        \begin{tikzpicture}[
                every node/.style={circle, draw, fill=black, inner sep=1.5pt},
            ]
            % Grafo 1 (Q3)
            \begin{scope}[shift={(0,0)}]
                \node[label=below left:4]  (a1) at (0,0) {};
                \node[label=below right:3] (a2) at (2,0) {};
                \node[label=above right:2] (a3) at (2,2) {};
                \node[label=above left:1] (a4) at (0,2) {};
                \node[label=above right:8] (b1) at (0.5,0.5) {};
                \node[label=above left:7] (b2) at (1.5,0.5) {};
                \node[label=above left:6] (b3) at (1.5,1.5) {};
                \node[label=above right:5] (b4) at (0.5,1.5) {};
                \draw[thick] (a1) -- (a2) -- (a3) -- (a4) -- (a1);
                \draw[thick] (b1) -- (b2) -- (b3) -- (b4) -- (b1);
                \draw[thick] (a1) -- (b1); \draw[thick] (a2) -- (b2); \draw[thick] (a3) -- (b3); \draw[thick] (a4) -- (b4);
            \end{scope}

            % Grafo 2 (Q3 con cruces)
            \begin{scope}[shift={(5,1)}]
                \node (a1) at (-1, 1) {};
                \node (a2) at (1, 1) {};
                \node (a3) at (1, -1) {};
                \node (a4) at (-1, -1) {};

                \node (b1) at (-0.5, 0.5) {};
                \node (b2) at (0.5, 0.5) {};
                \node (b3) at (0.5, -0.5) {};
                \node (b4) at (-0.5, -0.5) {};

                \draw[thick]
                (a1) -- (a2) -- (a3) -- (a4) -- (a1)
                (b1) -- (b2) -- (b3) -- (b4) -- (b1)
                (a1) -- (b1) (a1) -- (b4)
                (a2) -- (b2) (a2) -- (b3)
                (a3) -- (b3) (a3) -- (b2)
                (a4) -- (b4) (a4) -- (b1);
            \end{scope}
        \end{tikzpicture}
    \end{center}
\end{ejercicio}
\begin{proof}
    Notemos que si \( G_1 \) y \( G_2 \) son los dos grafos, luego \( G_1 \) y \( \overline{G_2} \) son dos \( Q_3 \) grafos. Es decir, tienen \( 8 \) vértices y \( 12 \) aristas cada uno y exactamente tres aristas inciden en cada vértice.

    Consideremos la función biyectiva \( f : V(G_1) \to V\left(\overline{G_2}\right) \) dada por: \[
        f = \begin{cases}
            1 & \mapsto 6, \quad 2 \mapsto 1, \\
            3 & \mapsto 7, \quad 4 \mapsto 4, \\
            5 & \mapsto 8, \quad 6 \mapsto 3, \\
            7 & \mapsto 2, \quad 8 \mapsto 5  \\
        \end{cases}
    \]

    Aplicando esta función, obtenemos:  \[
        \begin{cases}
            12 & \mapsto 61, \quad 23 \mapsto 17, \\
            34 & \mapsto 74, \quad 41 \mapsto 46, \\
            56 & \mapsto 83, \quad 67 \mapsto 32, \\
            78 & \mapsto 25, \quad 85 \mapsto 58, \\
            15 & \mapsto 68, \quad 26 \mapsto 13, \\
            37 & \mapsto 72, \quad 48 \mapsto 45.
        \end{cases}
    \]
    Obteniendo el isomorfismo entre las aristas de \( G_1 \) y \( \overline{G_2} \).
\end{proof}

\begin{ejercicio}{10}
    Determinar, para los grafos de abajo, qué pares de ellos son isomorfos y qué pares no lo son. En caso afirmativo, hallar un isomorfismo, dándole previamente nombres a los vértices. En caso negativo, dar argumentos estructurales.
    \begin{center}
        \begin{tikzpicture}[
                every node/.style={
                        circle,
                        fill=black,
                        inner sep=1.5pt,
                    },
                every path/.style={
                        draw=black,
                        line width=0.7pt
                    }
            ]

            \begin{scope}
                \node[label=above:$a$] (a) at (0.80,2.00) {};
                \node[label=above:$b$] (b) at (0.00,1.60) {};
                \node[label=above:$c$] (c) at (1.60,1.60) {};
                \node[label=above right:$d$] (d) at (1.80,0.80) {};
                \node[label=below:$e$] (e) at (1.25,0.00) {};
                \node[label=below:$f$] (f) at (0.45,0.00) {};
                \node[label=above left:$g$] (g) at (-0.15,0.80) {};

                \draw (a)--(b); \draw (a)--(c); \draw (a)--(d); \draw (a)--(g);
                \draw (b)--(c); \draw (b)--(f); \draw (b)--(g);
                \draw (c)--(d); \draw (c)--(e);
                \draw (d)--(e); \draw (d)--(f);
                \draw (e)--(f); \draw (e)--(g);
                \draw (f)--(g);
            \end{scope}

            \begin{scope}[xshift=3cm]
                \node[label=above:$a$] (a) at (0.80,2.00) {};
                \node[label=above:$b$] (b) at (0.00,1.60) {};
                \node[label=above:$c$] (c) at (1.60,1.60) {};
                \node[label=above right:$d$] (d) at (1.80,0.80) {};
                \node[label=below:$e$] (e) at (1.25,0.00) {};
                \node[label=below:$f$] (f) at (0.45,0.00) {};
                \node[label=above left:$g$] (g) at (-0.15,0.80) {};

                % Aristas
                \draw (a)--(b); \draw (a)--(d); \draw (a)--(c); \draw (a)--(g);
                \draw (b)--(f); \draw (b)--(d); \draw (b)--(g);
                \draw (c)--(d); \draw (c)--(e); \draw (c)--(f);
                \draw (d)--(e);
                \draw (e)--(f); \draw (e)--(g);
                \draw (f)--(g);
            \end{scope}

            \begin{scope}[xshift=6cm]
                \node[label=above:$a$] (a) at (0.80, 2.00) {};
                \node[label=above:$b$] (b) at (0.00, 1.60) {};
                \node[label=above:$c$] (c) at (1.60, 1.60) {};
                \node[label=above left:$d$] (d) at (-0.20, 0.75) {};
                \node[label=above right:$e$] (e) at (1.80, 0.75) {};
                \node[label=below:$f$] (f) at (0.30, 0.00) {};
                \node[label=below:$g$] (g) at (1.30, 0.00) {};
                % Aristas exteriores
                \draw (a)--(b);
                \draw (b)--(d);
                \draw (d)--(f);
                \draw (f)--(g);
                \draw (g)--(e);
                \draw (e)--(c);
                \draw (c)--(a);

                % Aristas interiores
                \draw (b)--(c);
                \draw (b)--(e);
                \draw (c)--(d);
                \draw (d)--(g);
                \draw (e)--(f);
                \draw (a)--(f);
                \draw (a)--(g);
            \end{scope}

            \begin{scope}[xshift=9cm]
                \node[label=above left:$a$] (a) at (0.43, 1.64) {};
                \node[label=above right:$b$] (b) at (1.28, 1.64) {};
                \node[label=right:$c$] (c) at (1.70, 0.90) {};
                \node[label=below right:$d$] (d) at (1.28, 0.16) {};
                \node[label=below left:$e$] (e) at (0.43, 0.16) {};
                \node[label=left:$f$] (f) at (0.00, 0.90) {};

                \node[label=above:$g$] (g) at (0.85, 0.90) {};

                \draw (a)--(b); \draw (b)--(c); \draw (c)--(d);
                \draw (d)--(e); \draw (e)--(f); \draw (f)--(a);

                \draw (a)--(c); \draw (c)--(e);
                \draw (d)--(f); \draw (f)--(b);

                \draw (a)--(g); \draw (b)--(g);
                \draw (d)--(g); \draw (e)--(g);
            \end{scope}

            \begin{scope}[xshift=12cm]
                \node[label=above:$a$] (a) at (1.00, 1.90) {};
                \node[label=above right:$b$] (b) at (1.78, 1.52) {};
                \node[label=right:$c$] (c) at (1.98, 0.68) {};
                \node[label=below right:$d$] (d) at (1.43, 0.00) {};
                \node[label=below left:$e$] (e) at (0.57, 0.00) {};
                \node[label=left:$f$] (f) at (0.02, 0.68) {};
                \node[label=above left:$g$] (g) at (0.22, 1.52) {};

                % 1. Aristas del perímetro
                \draw (a)--(b); \draw (b)--(c); \draw (c)--(d);
                \draw (d)--(e); \draw (e)--(f); \draw (f)--(g); \draw (g)--(a);

                % 3. Diagonales largas (salto de 2 nodos)
                \draw (a)--(d); \draw (b)--(e); \draw (c)--(f);
                \draw (d)--(g); \draw (e)--(a); \draw (f)--(b); \draw (g)--(c);
            \end{scope}
        \end{tikzpicture}
    \end{center}
\end{ejercicio}
\begin{proof}
    Enumeremos los grafos de izquierda a derecha como \( G_1, \ldots, G_5 \).
    Veamos si \( G_1 \) es isomorfo a \( G_2 \). Para ello, podemos intentar encontrar una correspondencia entre los vértices de ambos grafos que preserve las aristas.
    Notemos que: \begin{center}
        \begin{tikzpicture}[
                every node/.style={
                        circle,
                        fill=black,
                        inner sep=1.5pt,
                    },
                every path/.style={
                        draw=black,
                        line width=0.7pt
                    }
            ]
            \node[label=above:$a$] (a) at (0.80,2.00) {};
            \node[label=above:$b$] (b) at (0.00,1.60) {};
            \node[label=above:$d$] (d) at (1.60,1.60) {};
            \node[label=above right:$c$] (c) at (1.80,0.80) {};
            \node[label=below:$e$] (e) at (1.25,0.00) {};
            \node[label=below:$f$] (f) at (0.45,0.00) {};
            \node[label=above left:$g$] (g) at (-0.15,0.80) {};

            % Aristas
            \draw (a)--(b);
            \draw (a)--(d);
            \draw (a)--(c);
            \draw (a)--(g);

            \draw (b)--(f);
            \draw (b)--(d);
            \draw (b)--(g);

            \draw (c)--(d);
            \draw (c)--(e);
            \draw (c)--(f);

            \draw (d)--(e);

            \draw (e)--(f);
            \draw (e)--(g);

            \draw (f)--(g);
        \end{tikzpicture}
    \end{center}
    es el grafo obtenido via el isomorfismo que permuta \( c \) por \( d \) y deja todo lo demás igual.

    Veamos que \( G_3 \) no es isomorfo a \( G_1 \) y, por lo tanto, a \( G_2 \) tampoco. Para ello contemos la cantidad de triángulos en cada grafo. Si los grafos son isomorfos, deben tener la misma cantidad de triángulos.
    El grafo \( G_1 \) tiene los siguientes triángulos: \[
        \{ a,b,c \}, \,
        \{ a,c,d \}, \,
        \{ a,b,g \}, \,
        \{ g,b,f \}, \,
        \{ g,f,e \}, \,
        \{ f,d,e \}, \,
        \{ e,c,d \}
    \]
    Mientras que el grafo \( G_3 \) tiene: \[
        \{ a,b,c \}, \,
        \{ a,f,g \}, \,
        \{ b,c,e \}, \,
        \{ c,b,d \}, \,
        \{ e,f,g \}, \,
        \{ d,f,g \}
    \]
    Por lo tanto \( G_3 \) no es isomorfo a \( G_1 \) ni a \( G_2 \).
    Si contamos la cantidad de triángulos que posee \( G_4 \) llegamos a la conclusión de que tiene \( 6 \), por esto \( G_4 \) no es isomorfo a \( G_1 \) ni a \( G_2 \), veamos si lo es a \( G_3 \).

    Para ver que \( G_4 \) es isomorfo a \( G_3 \) notemos que podemos permutar de la siguiente forma:\[
        a \mapsto c, \,
        b \mapsto b, \,
        c \mapsto a, \,
        d \mapsto g, \,
        e \mapsto f, \,
        f \mapsto e, \,
        g \mapsto d
    \]
    obteniendo así un isomorfismo entre \( G_3 \) y \( G_4 \).

    Finalmente, veamos que \( G_5 \) no es isomorfo ni a \( G_3 \) ni a \( G_4 \), para ello, notemos que la cantidad de triángulos que posee \( G_5 \) es \( 7 \): \[
        \{ a,b,e \}, \, \{b,c,f\}, \, \{c,d,g\}, \, \{d,e,a\}, \, \{e,f,b\}, \, \{f,g,c\}, \, \{g,a,d\}
    \]
    Veamos si es isomorfo a \( G_1 \), para ello estudiemos los complementos de \( G_1 \) y \( G_5 \) respectivamente:
    \begin{center}
        \begin{tikzpicture}[
                every node/.style={
                        circle,
                        fill=black,
                        inner sep=1.5pt,
                    },
                every path/.style={
                        draw=black,
                        line width=0.7pt
                    }
            ]

            \begin{scope}
                \node[label=above:$a$] (a) at (0.80,2.00) {};
                \node[label=above:$b$] (b) at (0.00,1.60) {};
                \node[label=above:$c$] (c) at (1.60,1.60) {};
                \node[label=above right:$d$] (d) at (1.80,0.80) {};
                \node[label=below:$e$] (e) at (1.25,0.00) {};
                \node[label=below:$f$] (f) at (0.45,0.00) {};
                \node[label=above left:$g$] (g) at (-0.15,0.80) {};

                \draw (a)--(e); \draw (a)--(f);
                \draw (b)--(d); \draw (b)--(e);
                \draw (c)--(f); \draw (c)--(g);
                \draw (d)--(g);
            \end{scope}

            \begin{scope}[xshift=3.6cm]
                \node[label=above:$a$] (a) at (1.00, 1.90) {};
                \node[label=above right:$b$] (b) at (1.78, 1.52) {};
                \node[label=right:$c$] (c) at (1.98, 0.68) {};
                \node[label=below right:$d$] (d) at (1.43, 0.00) {};
                \node[label=below left:$e$] (e) at (0.57, 0.00) {};
                \node[label=left:$f$] (f) at (0.02, 0.68) {};
                \node[label=above left:$g$] (g) at (0.22, 1.52) {};

                \draw (a)--(c); \draw (b)--(d); \draw (c)--(e);
                \draw (d)--(f); \draw (e)--(g); \draw (f)--(a);
                \draw (g)--(b);
            \end{scope}
        \end{tikzpicture}
    \end{center}
    Nos basta con notar que ambos complementos se pueden llevar al mismo grafo de 7 vértices, lo que demuestra que \( \overline{G_1} \) es isomorfo a \( \overline{G_5} \) y, por el \emph{Ejercicio 8}, se concluye que \( G_1 \) es isomorfo a \( G_5 \).
\end{proof}

\vspace{0.5cm}

\begin{definition}\label{def:petersen}
    Un grafo de Petersen es un grafo simple cuyos vértices son subconjuntos de dos elementos de un conjunto de cinco elementos tales que dos vértices son adyacentes si y sólo si los subconjuntos correspondientes son disjuntos.
\end{definition}
Un ejemplo del grafo de Petersen se puede construir tomando el conjunto de cinco elementos \( \{1,2,3,4,5\} \) y considerando todos los subconjuntos de dos elementos como vértices. Dos vértices son adyacentes si y sólo si los subconjuntos correspondientes son disjuntos.
\begin{center}
    \begin{tikzpicture}[
            vertex/.style={circle, draw, fill=black, inner sep=1.5pt, outer sep=0pt},
        ]

        \def\R{2.8}
        \def\r{1.3}

        \node[vertex] (o1) at (90:\R) {};
        \node[above=3pt] at (o1) {12};

        \node[vertex] (o2) at (18:\R) {};
        \node[above right=2pt] at (o2) {34};

        \node[vertex] (o3) at (306:\R) {};
        \node[below right=2pt] at (o3) {51};

        \node[vertex] (o4) at (234:\R) {};
        \node[below left=2pt] at (o4) {23};

        \node[vertex] (o5) at (162:\R) {};
        \node[above left=2pt] at (o5) {45};

        \node[vertex] (i1) at (90:\r) {};
        \node[right=3pt] at (i1) {35};

        \node[vertex] (i2) at (18:\r) {};
        \node[below right=2pt] at (i2) {52};

        \node[vertex] (i3) at (306:\r) {};
        \node[left=3pt] at (i3) {24};

        \node[vertex] (i4) at (234:\r) {};
        \node[below right] at (i4) {41};

        \node[vertex] (i5) at (162:\r) {};
        \node[above right=2pt] at (i5) {13};

        \draw[thick] (o1) -- (o2) -- (o3) -- (o4) -- (o5) -- (o1);

        \draw[thick] (i1) -- (i3) -- (i5) -- (i2) -- (i4) -- (i1);

        \foreach \x in {1,...,5} {
                \draw[thick] (o\x) -- (i\x);
            }

    \end{tikzpicture}
\end{center}

\begin{ejercicio}{11}
    Probar que los grafos de abajo son isomorfos al grafo de Petersen. Hacerlo asociándole a los vértices conjuntos de dos elementos de manera tal que la regla de adyacencia entre vértices sea la misma que en la definición del grafo de Petersen.
    \begin{center}
        \begin{tikzpicture}[
                scale=.75,
                vertex/.style={circle, draw, fill=black, inner sep=1.5pt, outer sep=0pt}
            ]

            \begin{scope}
                \def\R{2.5}
                \def\r{0.8}

                \node[vertex, label=above:$15$] (v1) at (120:\R) {};
                \node[vertex, label=above:$24$] (v2) at (60:\R) {};
                \node[vertex, label=above right:$13$] (v3) at (0:\R) {};
                \node[vertex, label=below right:$25$] (v4) at (300:\R) {};
                \node[vertex, label=below left:$14$] (v5) at (240:\R) {};
                \node[vertex, label=above left:$23$] (v6) at (180:\R) {};

                \node[vertex, label=above:$12$] (c)   at (0,0) {};       % Centro
                \node[vertex, label=above left: $35$] (tli) at (150:\r) {};    % Arriba-Izquierda interior
                \node[vertex, label=above right: $34$] (tri) at (30:\r) {};     % Arriba-Derecha interior
                \node[vertex, label=below: $45$] (bi)  at (270:\r) {};    % Abajo interior

                \draw[thick] (v1) -- (v2) -- (v3) -- (v4) -- (v5) -- (v6) -- (v1);

                \draw[thick] (c) -- (tli);
                \draw[thick] (c) -- (tri);
                \draw[thick] (c) -- (bi);

                \draw[thick] (tli) -- (v2);
                \draw[thick] (tli) -- (v5);

                \draw[thick] (tri) -- (v1);
                \draw[thick] (tri) -- (v4);

                \draw[thick] (bi) -- (v6);
                \draw[thick] (bi) -- (v3);
            \end{scope}
            \begin{scope}[xshift=7cm, scale=1.5, yshift=-.5cm]

                % Definir coordenadas clave
                \def\radius{2.0} % Radio para los vértices exteriores

                % Vértices exteriores (un triángulo equilátero)
                \coordinate (Top) at (90:\radius);
                \coordinate (Left) at (210:\radius);
                \coordinate (Right) at (330:\radius);

                % Puntos medios de los lados exteriores
                \coordinate (MLeft)  at ($(Top)!0.5!(Left)$);
                \coordinate (MRight) at ($(Top)!0.5!(Right)$);
                \coordinate (MBottom) at ($(Left)!0.5!(Right)$);

                % Vértices interiores
                \coordinate (ITop)   at (90:0.8);  % Arriba en el eje vertical
                \coordinate (ILeft)  at (210:0.8); % Izquierda
                \coordinate (IRight) at (330:0.8); % Derecha
                \coordinate (ICenter) at (0,0);     % Centro

                % Punto de intersección entre la base exterior y la línea central
                \coordinate (BaseCenter) at ($(MBottom)!1!(ICenter)$);

                % --- Dibujar nodos (vértices) ---
                \node[vertex, label=above:$15$] (v1) at (Top)   {};
                \node[vertex, label=below left:$13$] (v2) at (Left)  {};
                \node[vertex, label=below right:$14$] (v3) at (Right) {};

                \node[vertex, label=above left: $24$] (v4) at (MLeft) {};
                \node[vertex, label=above right: $23$] (v5) at (MRight) {};
                \node[vertex, label=below: $25$] (v6) at (MBottom) {};

                \node[vertex, label=above left: $34$] (v7) at (ITop) {};
                \node[vertex, label=above: $45$] (v8) at (ILeft) {};
                \node[vertex, label=below: $35$] (v9) at (IRight) {};

                \node[vertex, label=above left: $12$] (v10) at (ICenter) {};

                % --- Dibujar aristas ---
                % Lados del triángulo exterior
                \draw[thick] (v1) -- (v2);
                \draw[thick] (v1) -- (v3);
                \draw[thick] (v2) -- (v3);

                % Líneas rectas interiores radiales
                \draw[thick] (v1) -- (v10);
                \draw[thick] (v2) -- (v10);
                \draw[thick] (v3) -- (v10);

                \draw[thick] (v6) to[in=-30] (v7);
                \draw[thick] (v5) to[out=-15, in=-15] (v8);
                \draw[thick] (v4) to[out=15, in=15] (v9);
            \end{scope}

            \begin{scope}[xshift=14cm]
                \def\R{2.5}

                \node[vertex, label=above:12] (v0) at (90:\R) {};   % Arriba
                \node[vertex, label=above right:35] (v7) at (45:\R) {};   % Arriba-Derecha
                \node[vertex, label=right:24] (v6) at (0:\R) {};    % Medio-Derecha
                \node[vertex, label=below right:13] (v5) at (315:\R) {};  % Abajo-Derecha
                \node[vertex, label=below:45] (v4) at (270:\R) {};  % Abajo
                \node[vertex, label=below left:23] (v3) at (225:\R) {};  % Abajo-Izquierda
                \node[vertex, label=left:15] (v2) at (180:\R) {};  % Medio-Izquierda
                \node[vertex, label=above left:34] (v1) at (135:\R) {};  % Arriba-Izquierda

                \node[vertex, label=below left:25] (i1) at (-0.9, -0.65) {}; % Interior Izquierdo
                \node[vertex, label=below right:14] (i2) at (0.9, -0.65) {};  % Interior Derecho

                \draw[thick] (v0) -- (v1) -- (v2) -- (v3) -- (v4) -- (v5) -- (v6) -- (v7) -- (v0);

                \draw[thick] (v0) -- (v4); % Línea vertical central
                \draw[thick] (v2) -- (v6); % Línea horizontal central

                \draw[thick] (i1) -- (i2);

                \draw[thick] (i1) -- (v1);
                \draw[thick] (i2) -- (v7);

                \draw[thick] (i1) -- (v5);
                \draw[thick] (i2) -- (v3);
            \end{scope}
        \end{tikzpicture}
    \end{center}
    \begin{center}
        \begin{tikzpicture}[
                scale=.75,
                vertex/.style={circle, draw, fill=black, inner sep=1.5pt, outer sep=0pt}
            ]
            \def\R{2.5}

            \foreach \i in {0,...,8} {
                    \node[vertex] (v\i) at (90 - \i * 40:\R) {};
                }

            \node[vertex, label={above:34}]        (v0) at (90:\R) {};   % Arriba
            \node[vertex, label={above right:15}]  (v1) at (50:\R) {};   % Arriba-Derecha
            \node[vertex, label={right:23}]        (v2) at (10:\R) {};   % Derecha-Arriba
            \node[vertex, label={below right:45}]  (v3) at (330:\R) {};  % Derecha-Abajo
            \node[vertex, label={below:13}]        (v4) at (290:\R) {};  % Abajo-Derecha
            \node[vertex, label={below:24}]        (v5) at (250:\R) {};  % Abajo-Izquierda
            \node[vertex, label={below left:35}]   (v6) at (210:\R) {};  % Izquierda-Abajo
            \node[vertex, label={left:14}]         (v7) at (170:\R) {};  % Izquierda-Arriba
            \node[vertex, label={above left:25}]   (v8) at (130:\R) {};  % Arriba-Izquierda

            \node[vertex, label={below:12}] (c) at (0,0) {};

            \draw[thick] (v0) -- (v1) -- (v2) -- (v3) -- (v4) -- (v5) -- (v6) -- (v7) -- (v8) -- (v0);

            \draw[thick] (c) -- (v0); % Hacia el nodo superior
            \draw[thick] (c) -- (v3); % Hacia el nodo inferior-derecho
            \draw[thick] (c) -- (v6); % Hacia el nodo inferior-izquierdo

            \draw[thick] (v7) -- (v2); % Línea horizontal que pasa apenas por encima del centro
            \draw[thick] (v8) -- (v4); % Línea que cruza desde arriba-izquierda hacia abajo-derecha
            \draw[thick] (v1) -- (v5); % Línea que cruza desde arriba-derecha hacia abajo-izquierda
        \end{tikzpicture}
    \end{center}
\end{ejercicio}

\clearpage

\section{Práctica II}
\begin{definition}
    Un \textbf{recorrido} (o paseo) de longitud \( k \) en un grafo \( G \) es una sucesión alternada de vértices y aristas \( v_0, e_1, v_1, e_2, \ldots, e_k, v_k \) tal que, para todo \( i \in \{1, \ldots, k\} \), los extremos de \( e_i \) son \( v_{i-1} \) y \( v_i \). El recorrido se dice \textbf{cerrado} si \( v_0 = v_k \).
\end{definition}

\begin{definition}
    Una \textbf{cadena} es un recorrido en el que no se repite ninguna arista.
\end{definition}

\begin{definition}
    Un \textbf{camino} es un recorrido en el que no se repite ningún vértice (y, por lo tanto, tampoco ninguna arista).
\end{definition}

\begin{definition}
    Un \textbf{ciclo} es una cadena cerrada no trivial (es decir, con al menos una arista) en la que no se repite ningún vértice, salvo el primero y el último, que coinciden.
\end{definition}

\begin{definition}\label{def:k4}
    Dado \( n \geq 3 \), el \textbf{grafo completo} \( K_n \) es el grafo simple de \( n \) vértices en el que cada par de vértices distintos es adyacente.
\end{definition}

\begin{center}
    \begin{tikzpicture}[every node/.style={circle, draw, fill=black, inner sep=1.5pt}]
        \node[label=above left:1] (v1) at (0,2) {};
        \node[label=above right:2] (v2) at (2,2) {};
        \node[label=below right:3] (v3) at (2,0) {};
        \node[label=below left:4] (v4) at (0,0) {};

        \draw[thick] (v1) -- (v2) -- (v3) -- (v4) -- (v1);
        \draw[thick] (v1) -- (v3);
        \draw[thick] (v2) -- (v4);
    \end{tikzpicture}
    \\[.25cm]
    El grafo completo \( K_4 \).
\end{center}

\begin{ejercicio}{1}
    Determinar si las siguientes afirmaciones son verdaderas o falsas. Justificar en cada caso:
    \begin{enumerate}
        \item \( K_4 \) contiene un recorrido que no es una cadena.
        \item \( K_4 \) contiene una cadena que no es cerrada y no es un camino.
        \item \( K_4 \) contiene una cadena cerrada no trivial (con al menos una arista) que no es un ciclo.
    \end{enumerate}
\end{ejercicio}
\begin{proof}
    Sea \( K_4 \) el grafo como en \textbf{Definición\,\ref{def:k4}}.\begin{enumerate}
        \item \( K_4 \) contiene un recorrido que no es una cadena, por ejemplo, un recorrido que repite una arista \( 12 - 21 \).
        \item \( K_4 \) contiene una cadena que no es cerrada y no es un camino, por ejemplo, una cadena que repite un vértice pero no una arista \( 12 - 23 - 31 - 14 \).
        \item Para que una cadena cerrada no trivial no sea un ciclo, la cadena debe repetir algún vértice (además del primero y el último). Si un recorrido pasa por un vértice, utiliza dos aristas distintas \( \implies \) Para que un vértice se repita en una cadena, debe tener grado al menos 4 dentro del subgrafo formado por el recorrido. Como en \( K_4 \) el grado máximo de cualquier vértice es 3 no puede existir ningún vértice con al menos 4 aristas incidentes y es imposible visitarlo más de una vez sin repetir aristas \( \therefore \) (c) es falsa.
    \end{enumerate}
\end{proof}

\clearpage

\begin{definition}
    El grado de un vértice \( v \) en un grafo \( G \), denotado \( \deg(v) \), es el número de aristas incidentes a \( v \).
\end{definition}

\begin{definition}
    La longitud de un recorrido, cadena o camino en un grafo es el número de aristas que contiene. Si \( R \) es el recorrido, notaremos \( \ell(R) \) a su longitud.
\end{definition}

\begin{definition}
    Un grafo de Petersen es un grafo simple cuyos vértices son subconjuntos de dos elementos de un conjunto de cinco elementos tales que dos vértices son adyacentes si y sólo si los subconjuntos correspondientes son disjuntos.
\end{definition}
Un ejemplo del grafo de Petersen se puede construir tomando el conjunto de cinco elementos \( P = \{1,2,3,4,5\} \) y considerando todos los subconjuntos de dos elementos como vértices. Dos vértices son adyacentes si y sólo si los subconjuntos correspondientes son disjuntos.
\begin{center}
    \begin{tikzpicture}[
            vertex/.style={circle, draw, fill=black, inner sep=1.5pt, outer sep=0pt},
        ]

        \def\R{2.8}
        \def\r{1.3}

        \node[vertex] (o1) at (90:\R) {};
        \node[above=3pt] at (o1) {12};

        \node[vertex] (o2) at (18:\R) {};
        \node[above right=2pt] at (o2) {34};

        \node[vertex] (o3) at (306:\R) {};
        \node[below right=2pt] at (o3) {51};

        \node[vertex] (o4) at (234:\R) {};
        \node[below left=2pt] at (o4) {23};

        \node[vertex] (o5) at (162:\R) {};
        \node[above left=2pt] at (o5) {45};

        \node[vertex] (i1) at (90:\r) {};
        \node[right=3pt] at (i1) {35};

        \node[vertex] (i2) at (18:\r) {};
        \node[below right=2pt] at (i2) {52};

        \node[vertex] (i3) at (306:\r) {};
        \node[left=3pt] at (i3) {24};

        \node[vertex] (i4) at (234:\r) {};
        \node[below right] at (i4) {41};

        \node[vertex] (i5) at (162:\r) {};
        \node[above right=2pt] at (i5) {13};

        \draw[thick] (o1) -- (o2) -- (o3) -- (o4) -- (o5) -- (o1);

        \draw[thick] (i1) -- (i3) -- (i5) -- (i2) -- (i4) -- (i1);

        \foreach \x in {1,...,5} {
                \draw[thick] (o\x) -- (i\x);
            }

    \end{tikzpicture}
\end{center}

\begin{ejercicioEntregar}{2}
    Probar que el grafo de Petersen no tiene ciclos de longitud 7.
\end{ejercicioEntregar}
\begin{proof}
    Es fácil ver que el grafo de Petersen no puede contener cíclos de longitud 3 pues por como está construido tener un ciclo de longitud 3 implicaría tener 3 subconjuntos dos a dos disjuntos de 2 elementos del conjunto \( P = \{1,2,3,4,5\} \), lo cual es imposible.

    Para ver que no puede haber un ciclo de longitud 4, supongamos que existe uno de la forma \( u - v - w - x - u \). Por la definición de adyacencia, los vértices consecutivos deben ser conjuntos disjuntos. Sin pérdida de generalidad, supongamos que \( u = \{1, 2\} \). Como \( v \) es adyacente a \( u \), debe ser disjunto, por lo que podemos asumir que \( v = \{3, 4\} \).

    El vértice \( w \) debe ser adyacente a \( v \), lo que implica que \( w \cap \{3, 4\} = \varnothing \). Por lo tanto, los dos elementos de \( w \) deben elegirse del conjunto restante \( \{1, 2, 5\} \). Además, \( w \neq u \) (pues son vértices distintos del ciclo), lo que nos obliga a que \( w \) contenga al 5 y a un elemento de \( u \). Supongamos s.p.d.g que \( w = \{1, 5\} \).

    Finalmente, el vértice \( x \) debe cerrar el ciclo, por lo que debe ser adyacente tanto a \( w = \{1, 5\} \) como a \( u = \{1, 2\} \). Esto exige que \( x \) sea disjunto de ambos conjuntos, lo que significa que \( x \) no puede contener a los elementos 1, 2 ni 5.

    Los únicos elementos disponibles en \( P \) para formar el subconjunto de dos elementos \( x \) son el 3 y el 4, forzando a que \( x = \{3, 4\} \). Pero esto implica que \( x = v \), lo cual es una contradicción porque un ciclo de longitud 4 requiere cuatro vértices distintos. Por lo tanto, no puede existir un ciclo de longitud 4.

    Sea \( G \) el grafo de Petersen. Notemos en primer lugar que \( G \) es tal que todo vértice tiene grado 3.

    Supongamos por el absurdo que existe un ciclo \( C \) de longitud 7 en \( G \). Sea \( S = V(G) \setminus V(C) \). Como el ciclo tiene 7 vértices y \( G \) tiene 10, debe ocurrir que \( |S| = 3 \).

    Dado que \( \forall v \in V(G) \) vale que \( \deg(v) = 3 \), cada vértice \( v \in V(C) \) utiliza dos de sus tres aristas incidentes para formar el ciclo \( C \). La arista restante de cada vértice de \( C \) no puede conectar con otro vértice de \( C \), ya que esto dividiría al ciclo en ciclos más pequeños de longitud 3 o 4 lo cual no es posible. Por lo tanto, esas 7 aristas deben conectar necesariamente con vértices de \( S \).

    Entonces tenemos 7 aristas que van desde \( C \) hacia los 3 vértices de \( S \). Sabemos que la suma de los grados de los vértices de \( S \) es \( 3 \times 3 = 9 \). Definamos \( k \) como el número de aristas internas en \( S \) (las que conectan dos vértices de \( S \)). Cada una de estas aristas internas consume 2 aristas del total disponible en el conjunto. Es decir, la cantidad de aristas que van de \( S \) a \( C \) está dada por \( 9 - 2k \). Igualando, tenemos \( 9 - 2k = 7 \implies 2k = 2 \implies k = 1 \).

    Esto implica que en \( S \) hay exactamente una arista interna, por lo que uno de los vértices de \( S \) es aislado dentro de ese subconjunto. Sea \( z \) este vértice aislado en \( S \). Sabemos que las 3 aristas de \( z \) van directamente a parar a \( C \), conectando a \( z \) con 3 vértices distintos del ciclo.

    Si enumeramos \( v_1, v_2, \ldots, v_7 \) a los vértices del ciclo en orden, entonces \( z \) debe estar conectado a tres de ellos, digamos \( v_i, v_j, v_k \). Sin pérdida de generalidad supongamos \( 1 \leq i < j < k \leq 7 \). Estas conexiones dividen al ciclo en 3 caminos sobre su contorno, cuyas longitudes son \( a = j-i \), \( b = k-j \), y \( c = 7 - k + i \).

    La suma de las longitudes de estos caminos es \( a + b + c = 7 \). Como los tres son enteros positivos, si todos fueran mayores o iguales a 3, su suma sería al menos 9. Por lo tanto, obligatoriamente alguno de estos caminos debe tener longitud 1 o 2.
    \begin{itemize}
        \item Si un camino tiene longitud 1, los dos vértices en \( C \) son adyacentes y, junto con \( z \), forman un ciclo de longitud 3.
        \item Si un camino tiene longitud 2, los vértices en \( C \) están separados por un solo vértice intermedio y, junto con \( z \), forman un ciclo de longitud 4.
    \end{itemize}
    Ambas situaciones son imposibles en el grafo de Petersen. Por lo tanto, no puede existir un ciclo de longitud 7.
\end{proof}

\vspace{0.5cm}

\begin{ejercicioEntregar}{3}
    Demostrar que, si todos los vértices de un grafo \( G \) simple y conexo tienen grado 2, entonces \( G \) es un ciclo. ¿Vale la misma conclusión si el grafo no es simple?
\end{ejercicioEntregar}
\begin{proof}
    Consideremos un camino maximal en \(G\), dado por la sucesión de vértices \(P = v_1, v_2, \ldots, v_k\).
    Analicemos el último vértice del camino, \(v_k\). Como \(\deg(v_k) = 2\), tiene exactamente dos aristas incidentes. Una de ellas lo conecta con \(v_{k-1}\). La otra arista debe conectarlo necesariamente con algún vértice que ya pertenece al camino \(P\). Si lo conectara con un vértice \(w \notin V(P)\), el camino podría extenderse a \(v_1, \ldots, v_k, w\), lo cual contradice que \(P\) sea un camino maximal.

    Sea \(v_j\) (con \(1 \leq j \leq k-2\)) ese vértice de \(P\) adyacente a \(v_k\). La unión de esta arista con la porción del camino forma un ciclo \(C = v_j, v_{j+1}, \ldots, v_k, v_j\).

    Ahora veamos que este ciclo abarca todo el grafo. Cada vértice en el subgrafo \(C\) ya tiene grado 2 utilizando únicamente las aristas que pertenecen al ciclo (una de entrada desde el vértice anterior y otra de salida hacia el siguiente). Como la hipótesis general indica que el grado de todo vértice en \(G\) es exactamente 2, ningún vértice de \(C\) puede tener aristas incidentes adicionales que lo conecten con vértices fuera de \(C\).

    Dado que \(G\) es un grafo conexo, es imposible que existan vértices de \(G\) desconectados de \(C\). Por lo tanto, no hay vértices fuera del ciclo, de modo que \(V(G) = V(C)\) y \(E(G) = E(C)\). En consecuencia, \(G\) es exactamente el ciclo \(C\).

    Si el grafo no es simple y contiene un lazo, entonces ese lazo por sí solo ya cumple con la condición de que el vértice involucrado tiene grado 2, pero como también es conexo, el grafo solo puede consistir en ese lazo, por lo tanto se concluye que el grafo es un ciclo de longitud 1. Un razonamiento análogo se puede aplicar a si contiene una arista múltiple.
\end{proof}

\begin{algorithm}
    \caption{Bubble Sort}
    \label{alg:bubblesort}
    \begin{algorithmic}[1] % El [1] es para enumerar cada línea
        \Procedure{BubbleSort}{$A, n$}
        \State \textbf{Entrada:} Una lista $A$ de $n$ elementos.
        \State \textbf{Salida:} La lista $A$ ordenada de forma ascendente.
        \vspace{0.2cm}
        \For{$i \gets 1$ \textbf{hasta} $n-1$}
        \For{$j \gets 1$ \textbf{hasta} $n-i$}
        \If{$A[j] > A[j+1]$}
        \State Intercambiar $A[j]$ y $A[j+1]$
        \EndIf
        \EndFor
        \EndFor
        \EndProcedure
    \end{algorithmic}
\end{algorithm}

\begin{ejercicio}{4}
    Sea \( G_n \) el grafo cuyos vértices son las permutaciones de \( \{1, 2, \ldots, n\} \), con dos permutaciones adyacentes si y sólo si difieren en un intercambio de un par de elementos consecutivos. Probar que \( G_n \) es conexo.
\end{ejercicio}
\begin{proof}
    La demostración es sencilla pues si demostramos que cualquier vértice del grafo está conectado al vértice correspondiente a la permutación identidad \( (1, 2, \ldots, n) \), entonces todos los vértices estarán conectados entre sí, lo que implica que el grafo es conexo.

    Sea cualquier permutación \( \sigma = (\sigma(1), \sigma(2), \ldots, \sigma(n)) \) de \( \{1, 2, \ldots, n\} \). Queremos mostrar que existe un camino en \( G_n \) que conecta \( \sigma \) con la permutación identidad \( (1, 2, \ldots, n) \).  Pero esto es muy sencillo pues si aplicamos el algoritmo Bubble Sort a \( \sigma \), cada intercambio de elementos consecutivos corresponde a una arista en \( G_n \), y al finalizar el algoritmo obtenemos la permutación identidad. Por lo tanto, existe un camino en \( G_n \) desde cualquier vértice hasta la permutación identidad, lo que demuestra que \( G_n \) es conexo.
\end{proof}

\vspace{0.5cm}

\begin{definition}
    Una \textbf{componente conexa} de un grafo \( G \) es un subgrafo conexo maximal de \( G \), es decir, un subgrafo conexo que no está contenido en ningún subgrafo conexo más grande de \( G \).
\end{definition}

\begin{ejercicioEntregar}{5}
    Sea \( G \) el grafo simple con conjunto de vértices \( \{1, 2, \ldots, 15\} \) tal que \( i \) y \( j \) son adyacentes si y sólo si su máximo común divisor es mayor que 1. Contar las componentes conexas de \( G \) y hallar un camino de longitud máxima en \( G \).
\end{ejercicioEntregar}
\begin{proof}
    Analicemos si \( G \) posee vértices aislados, es decir, vértices que no son adyacentes a ningún otro vértice. Un vértice \( i \) será aislado si y sólo si \(\gcd(i,j) = 1\) para todo \( j \neq i \). Observemos que los números primos mayores que 7 (es decir, 11, 13) serán aislados, ya que no tienen divisores comunes con ningún otro vértice en el conjunto \( \{1, 2, \ldots, 15\} \). Por lo tanto, los vértices aislados son \( 11 \) y \( 13 \). Como \( \gcd(i, 1) = 1 \) para todo \( i \neq 1 \), el vértice 1 es aislado.

    Ahora consideremos los vértices no aislados: \( \{2, 3, 4, 5, 6, 7, 8, 9, 10, 12, 14, 15\} \). Podemos agruparlos según sus divisores primos:
    \begin{itemize}
        \item Múltiplos de 2: \(2, 4, 6, 8, 10, 12, 14\)
        \item Múltiplos de 3: \(3, 6, 9, 12, 15\)
        \item Múltiplos de 5: \(5, 10, 15\)
        \item Múltiplos de 7: \(7, 14\)
    \end{itemize}

    Como algunos vértices pertenecen a más de un grupo, estos vértices actúan como puentes que conectan los diferentes subgrupos. Por ejemplo, el vértice 6 conecta los múltiplos de 2 y 3, el vértice 10 conecta los múltiplos de 2 y 5, y el vértice 14 conecta los múltiplos de 2 y 7. Por lo tanto, todos los vértices no aislados están conectados entre sí, formando una única componente conexa.

    En conclusión, \( G \) tiene 4 componentes conexas: \( \{ 1 \} \), \( \{11\} \), \( \{13\} \) y \( \{2, 3, 4, 5, 6, 7, 8, 9, 10, 12, 14, 15\} \).

    Un camino de longitud máxima en la componente conexa grande puede ser \[ 7 - 14 - 2 - 4 - 8 - 10 - 5 - 15 - 3 - 9 - 6 - 12 \]
\end{proof}

\vspace{0.5cm}

\begin{ejercicioEntregar}{6}
    Probar que un grafo es conexo si y sólo si, para cada partición de sus vértices en dos conjuntos no vacíos, existe una arista con extremos en ambos conjuntos.
\end{ejercicioEntregar}
\begin{proof}
    (\( \Rightarrow \)) Supongamos que \( G \) es conexo y consideremos una partición cualquiera de sus vértices en dos conjuntos no vacíos \( A \) y \( B \). Tomemos un vértice \( u \in A \) y un vértice \( v \in B \). Como \( G \) es conexo, existe un camino que une \( u \) con \( v \).

    Al recorrer los vértices de este camino partiendo desde \( u \in A \) hasta llegar a \( v \in B \), tiene que existir un primer vértice en el camino que pertenezca a \( B \). Llamémoslo \( y \). El vértice inmediatamente anterior a \( y \) en el camino, llamémoslo \( x \), debe estar en \( A \) por ser \( y \) el primero en \( B \). La arista que une a \( x \) con \( y \) es la arista buscada, ya que tiene un extremo en \( A \) y otro en \( B \).

    (\( \Leftarrow \)) Supongamos que para cada partición de \( V(G) \) en dos conjuntos disjuntos y no vacíos, existe al menos una arista que los conecta. Queremos probar que \( G \) es conexo.

    Sea \( u \in V(G) \). Definamos al conjunto \( A \) como la componente conexa de \( u \), es decir, el conjunto de todos los vértices de \( G \) a los que se puede llegar desde \( u \) mediante algún camino (notemos que \( u \in A \), por lo que \( A \neq \varnothing \)).

    Supongamos por el absurdo que \( G \) no es conexo. Entonces existe al menos un vértice al que no se puede llegar desde \( u \), lo que implica que el conjunto \( B = V(G) \setminus A \) es no vacío. Tenemos entonces una partición de \( V(G) \) en dos conjuntos no vacíos \( A \) y \( B \).

    Por hipótesis, debe existir una arista con un extremo \( x \in A \) y otro extremo \( y \in B \). Sin embargo, como \( x \in A \), existe un camino desde \( u \) hasta \( x \). Si a ese camino le agregamos la arista \( x-y \), obtenemos un camino desde \( u \) hasta \( y \). Esto significa que \( y \) es alcanzable desde \( u \), por lo que \( y \in A \), lo cual contradice que \( y \in B \).

    Esta contradicción proviene de suponer que \( B \) era no vacío. Por lo tanto, \( B = \varnothing \), lo que implica que \( A = V(G) \). Como todos los vértices son alcanzables desde \( u \), el grafo es conexo.

    \vspace{0.1cm}

    Entiendo que la recíproca se podría demostrar de una manera más sencilla en grafos finitos eligiendo la partición \( V(G) = \{ u \} \cup V(G) \setminus \{ u \} \). Tomando el vértice que cumple la hipótesis \( w_1 \in V(G) \setminus \{ u \} \) y aplicando inductivamente el argumento a \( A = V(G) \setminus \{ u, w_1 \} \) y \( B = \{ u, w_1\} \) hasta agotar todos los vértices del grafo o dar con \( v \). Pero no vale en grafos con vértices infinitos pues podríamos no terminar de recorrer todos los vértices.
\end{proof}

\vspace{0.5cm}

\begin{ejercicio}{7}
    Sean \( P \) y \( Q \) caminos de máxima longitud en un grafo conexo. Probar que \( P \) y \( Q \) tienen un vértice en común.
\end{ejercicio}
\begin{proof}
    Supongamos por el absurdo que \( P \) y \( Q \) no tienen ningún vértice en común y ambos son caminos de máxima longitud \( k \) en el grafo conexo \( G \). Dado que \( G \) es conexo, debe existir un camino que conecte algún vértice de \( P \) con algún vértice de \( Q \).

    Sean \( u \in P \) y \( v \in Q \) los extremos del camino más corto que conecta a \( P \) y \( Q \). El vértice \( u \) divide a \( P \) en dos subcaminos; tomemos el que tenga mayor o igual longitud (el cual medirá al menos \( k/2 \)). Análogamente, el vértice \( v \) divide a \( Q \) en dos subcaminos; tomemos también el de mayor o igual longitud (que medirá al menos \( k/2 \)).

    Entonces, al unir el subcamino de mayor longitud de \( P \) que termina en \( u \), el camino que conecta \( u \) con \( v \) (cuya longitud es al menos 1), y el subcamino de mayor longitud de \( Q \) que empieza en \( v \), obtenemos un nuevo camino simple cuya longitud es al menos \( k/2 + 1 + k/2 = k + 1 \). Esto es estrictamente mayor que \( k \), lo cual contradice la maximalidad de las longitudes de \( P \) y \( Q \).
\end{proof}

\vspace{0.5cm}

\begin{definition}
    Dado un grafo \( G \), el \textbf{complemento} de \( G \), denotado \( \overline{G} \), es el grafo con el mismo conjunto de vértices que \( G \) en el que dos vértices distintos son adyacentes si y sólo si no lo son en \( G \).
\end{definition}

\begin{definition}
    Un \textbf{vértice de corte} de un grafo conexo \( G \) es un vértice \( v \) tal que \( G - v \) tiene más de una componente conexa que \( G \).
\end{definition}

\begin{ejercicio}{8}
    Sea \( v \) un vértice de corte de un grafo simple \( G \). Probar que \( \overline{G} - v \) es conexo.
\end{ejercicio}
\begin{proof}
    Como \( v \) es un vértice de corte de \( G \), el subgrafo \( G - v \) tiene al menos dos componentes conexas. Sean \( C_1, C_2, \dots, C_k \) los conjuntos de vértices de dichas componentes conexas, con \( k \ge 2 \).

    Sean \( x, y \in V(\overline{G} - v) \) dos vértices distintos cualesquiera. Como \( V(\overline{G} - v) = V(G) \setminus \{v\} \), sabemos que \( x \) e \( y \) son vértices de \( G - v \). Veamos que existe un camino entre ellos en \( \overline{G} - v \) analizando dos casos:

    \begin{enumerate}
        \item[I:] \( x \) e \( y \) pertenecen a distintas componentes conexas de \( G - v \).
              Supongamos sin pérdida de generalidad que \( x \in C_1 \) e \( y \in C_2 \). Al pertenecer a componentes conexas distintas de \( G - v \), no existe ninguna arista que los una en \( G - v \), y por lo tanto tampoco existe en \( G \). Por la definición de grafo complemento, la arista \( xy \) debe existir en \( \overline{G} \). Como ni \( x \) ni \( y \) son el vértice \( v \), la arista \( xy \) se preserva en \( \overline{G} - v \). Así, \( x \) e \( y \) son adyacentes y están directamente conectados.
        \item[II:] \( x \) e \( y \) pertenecen a la misma componente conexa de \( G - v \).
              Supongamos que \( x, y \in C_1 \). Como \( k \ge 2 \), existe al menos otra componente conexa \( C_2 \) en \( G - v \). Tomemos un vértice cualquiera \( z \in C_2 \).
              Aplicando el razonamiento del anterior, como \( x \in C_1 \) y \( z \in C_2 \), la arista \( xz \) existe en \( \overline{G} - v \). De manera análoga, como \( y \in C_1 \) y \( z \in C_2 \), la arista \( yz \) también existe en \( \overline{G} - v \).
              Por lo tanto, existe el camino \( x - z - y \) en \( \overline{G} - v \), lo que demuestra que \( x \) e \( y \) están conectados.
    \end{enumerate}

    En ambos casos hemos encontrado un camino que conecta a \( x \) e \( y \) en \( \overline{G} - v \). Como \( x \) e \( y \) eran arbitrarios, vale que \( \overline{G} - v \) es conexo.

    Notar que el razonamiento es análogo al \textbf{Ejercicio 3} de la Práctica I.
\end{proof}
\vspace{0.5cm}

\begin{definition}
    Dos grafos simples \( G \) y \( H \) son \textbf{isomorfos} \( \iff \) existe una biyección \( f : V(G) \to V(H) \) tal que \( \{u,v\} \in E(G) \iff \{f(u), f(v)\} \in E(H) \) para todo \( u,v \in V(G) \).
\end{definition}

\begin{definition}
    Un grafo \( G \) se dice \textbf{auto-complementario} si es isomorfo a su complemento \( \overline{G} \).
\end{definition}

\begin{ejercicio}{9}
    Sea \( G \) un grafo auto-complementario. Probar que \( G \) tiene un vértice de corte si y sólo si \( G \) tiene un vértice de grado 1.
\end{ejercicio}
\begin{proof}
    Notar que para ser auto-complementario no trivial, \( G \) debe tener al menos 4 vértices.

    \( (\implies) \) Sea \( v \) un vértice de corte de \( G \). Entonces \( G - v \) es disconexo y sus vértices se particionan en las componentes conexas \( C_1, C_2, \dots, C_k \) con \( k \ge 2 \).
    En \( \overline{G} \), cualquier vértice de \( C_i \) es adyacente a todos los vértices de \( C_j \quad \forall i \neq j \).
    Debido a que el isomorfismo preserva la estructura, \( \overline{G} \) debe tener un vértice de corte. Sea \( w \) dicho vértice de corte en \( \overline{G} \).
    Si eliminamos \( w \) de \( \overline{G} \), los vértices restantes de \( V(G) \setminus \{v, w\} \) siguen formando un subgrafo conexo, ya entre los \( C_i \) asegura que exista un camino entre cualquier par de vértices (como \( n \ge 4 \), siempre hay suficientes vértices para mantener la conexión).
    Por lo tanto, la única forma en que la eliminación de \( w \) pueda desconectar a \( \overline{G} \) es si \( w \) aísla a \( v \) del resto del grafo. Esto implica que \( v \) está conectado únicamente a \( w \) en \( \overline{G} \), lo que significa que el grado de \( v \) en \( \overline{G} \) es 1. Al ser \( G \) y \( \overline{G} \) isomorfos, \( G \) también debe tener un vértice de grado 1.

    \( (\impliedby) \) Sea \( x \) un vértice de grado 1 en \( G \), y sea \( u \) su único vecino.
    Como \( G \) es auto-complementario no trivial, su cantidad de vértices es \( n \ge 4 \).
    Consideremos el subgrafo \( G - u \). Al eliminar el vértice \( u \), el vértice \( x \) pierde su única arista incidente, quedando aislado.
    Dado que originalmente el grafo tenía al menos 4 vértices, \( G - u \) contiene a \( x \) y al menos a otros dos vértices. Por lo tanto, \( G - u \) se descompone en al menos dos componentes conexas \( \therefore u \) es un vértice de corte en \( G \).
\end{proof}

\clearpage

\begin{definition}
    Dado un grafo \( G \) y un subconjunto de vértices \( S \subseteq V(G) \), el \textbf{subgrafo inducido} por \( S \), denotado \( G[S] \), es el grafo con conjunto de vértices \( S \) cuyas aristas son exactamente las aristas de \( G \) con ambos extremos en \( S \).
\end{definition}

\begin{definition}
    Un vértice \( v \) de un grafo \( G \) se dice \textbf{aislado} si \( \deg(v) = 0 \) i.e no tiene vecinos.
\end{definition}

\begin{ejercicioEntregar}{10}
    Sea \( G \) un grafo simple sin vértices aislados y con ningún subgrafo inducido que posea exactamente dos aristas. Probar que \( G \) es completo.
\end{ejercicioEntregar}

\begin{proof}
    Supongamos por el absurdo que \( G \) no es completo. Entonces, existen al menos dos vértices \( u, v \in V(G) \) tal que \( uv \notin E(G) \). Como \( G \) no tiene vértices aislados, \( \exists x \in V(G) \setminus \{ u, v\} \) adyacente a \( u \), y \( \exists y \in V(G) \setminus \{ u, v\} \) adyacente a \( v \).

    \textbf{Caso 1}: \( u \) y \( v \) tienen un vecino en común. \\
    Sea \( w \) un vértice adyacente tanto a \( u \) como a \( v \). Consideremos el subgrafo inducido por el conjunto de vértices \( S = \{u, v, w\} \). Las únicas aristas en este subgrafo son \( uw \) y \( vw \) (ya que \( uv \notin E(G) \)). Por lo tanto, \( G[S] \) tiene exactamente dos aristas, lo cual contradice la hipótesis.

    \textbf{Caso 2}: \( u \) y \( v \) no tienen vecinos en común. \\
    Como no comparten vecinos, el vértice \( x \) (vecino de \( u \)) no es adyacente a \( v \), y el vértice \( y \) (vecino de \( v \)) no es adyacente a \( u \). Notemos también que obligatoriamente \( x \neq y \). Consideremos si \( x \) e \( y \) son adyacentes entre sí:

    \begin{itemize}
        \item \textbf{Subcaso 2.1}: \( xy \notin E(G) \). \\
              Consideremos el subgrafo inducido por \( S = \{u, v, x, y\} \). Las únicas aristas presentes en este subgrafo son \( ux \) y \( vy \). Este subgrafo tiene exactamente dos aristas, lo cual es una contradicción.

        \item \textbf{Subcaso 2.2}: \( xy \in E(G) \). \\
              Consideremos el subgrafo inducido por \( S = \{u, x, y\} \). Sabemos que \( ux \in E(G) \) y \( xy \in E(G) \). Además, sabemos que \( uy \notin E(G) \) (porque \( u \) y \( v \) no comparten vecinos, e \( y \) es un vecino de \( v \)). Este subgrafo de tres vértices tiene exactamente dos aristas \( ux \) y \( xy \), contradiciendo la hipótesis.
    \end{itemize}

    En todos los casos posibles, la suposición de que \( G \) no es completo nos lleva a encontrar un subgrafo inducido con exactamente dos aristas \( \therefore G \) debe ser completo.
\end{proof}
\clearpage

\section{Práctica III}
Recordemos la noción de conjunto independiente (Definición~\ref{def:conjunto-independiente} de la Práctica~I).

\begin{definition}
    Un grafo \( G \) es \textbf{bipartito} si existe una partición de \( V(G) \) en dos conjuntos no vacíos \( A \) y \( B \) (a la que llamaremos una \textbf{bipartición} de \( G \)) tal que toda arista de \( G \) tiene un extremo en \( A \) y el otro en \( B \). Equivalentemente, \( A \) y \( B \) son conjuntos independientes de \( G \).
\end{definition}

\begin{ejercicio}{1}
    Determinar si el grafo de Petersen es bipartito y encontrar un conjunto independiente máximo en él, justificando por qué no se puede encontrar uno mayor.
\end{ejercicio}
\begin{proof}
    Propongo que el grafo de Petersen no es bipartito. Aplicaremos el algoritmo del \textbf{Ejercicio 11} de la Práctica I partiendo del vértice \( v = 12 \), recordando que dos vértices son adyacentes si y sólo si los subconjuntos correspondientes son disjuntos.

    \begin{center}
        \resizebox{\textwidth}{!}{%
            \begin{tabular}{cccc}
                \PetersenPaso{A,34,51,23,45}{35,52,24,41,13} &
                \PetersenPaso{A,B,51,23,B}{B,52,24,41,13}    &
                \PetersenPaso{A,B,51,A,B}{B,52,24,41,A}      &
                \PetersenPaso{A,B,B,A,B}{B,52,24,41,A}                                                                            \\
                \footnotesize Paso 1                         & \footnotesize Paso 2 & \footnotesize Paso 3 & \footnotesize Paso 4
            \end{tabular}%
        }
    \end{center}

    \begin{itemize}
        \item \textbf{Paso 1:} etiquetamos al vértice inicial \( 12 \) con \( A \).
        \item \textbf{Paso 2:} sus tres vecinos, \( 34 \), \( 45 \) y \( 35 \) (los únicos subconjuntos disjuntos de \( \{1,2\} \) entre los adyacentes a \( 12 \)), reciben el signo opuesto \( B \).
        \item \textbf{Paso 3:} procesamos el vecino \( 45 \); sus vecinos aún no etiquetados, \( 23 \) y \( 13 \), reciben el signo opuesto \( A \).
        \item \textbf{Paso 4:} procesamos el vecino \( 23 \) (con signo \( A \)); su vecino \( 51 \) recibe el signo opuesto \( B \). Pero \( 51 \) es también vecino de \( 34 \) (pues \( \{5,1\} \cap \{3,4\} = \varnothing \)), que ya tenía asignado el signo \( B \).
    \end{itemize}

    Encontramos así dos vértices adyacentes, \( 34 \) y \( 51 \), a los que el algoritmo asignó el mismo signo. Por lo que podemos concluir que el grafo de Petersen no es bipartito.

    \vspace{0.3cm}

    Encontremos ahora un conjunto independiente máximo. Consideremos
    \[
        S = \{12, 13, 14, 15\},
    \]
    es decir, todos los subconjuntos de dos elementos que contienen al \( 1 \). Es claro que dos elementos de \( S \) no son adyacentes en el grafo de Petersen. Así, \( S \) es un conjunto independiente de tamaño \( 4 \).

    Veamos que no existe un conjunto independiente de tamaño \( 5 \). Como los vértices del grafo de Petersen son los subconjuntos de dos elementos de \( \{1,\ldots,5\} \) y dos de ellos son no adyacentes si y sólo si se intersecan, un conjunto independiente se corresponde exactamente con una familia \( \mathcal{F} \) de subconjuntos de dos elementos de \( \{1,\ldots,5\} \) que se intersecan dos a dos. Pensemos a cada elemento de \( \mathcal{F} \) como una arista del grafo completo \( K_5 \) con vértices \( \{1,\ldots,5\} \); la condición dice entonces que dos aristas cualesquiera de \( \mathcal{F} \) comparten un extremo.

    \begin{itemize}
        \item Si todas las aristas de \( \mathcal{F} \) tienen un vértice \( x \) en común (\( \mathcal{F} \) es una \emph{estrella} centrada en \( x \)), entonces \( |\mathcal{F}| \leq 4 \), pues en \( K_5 \) hay exactamente 4 aristas incidentes a \( x \).
        \item Si no hay un vértice común a todas, tomemos \( e_1 = \{a,b\} \in \mathcal{F} \) y otra arista que la interseque, digamos \( e_2 = \{b,c\} \) con \( c \neq a \). Como no todas las aristas contienen a \( b \), existe \( e_3 \in \mathcal{F} \) con \( b \notin e_3 \). Para intersecar a \( e_1 \) sin usar \( b \), \( e_3 \) debe contener a \( a \); para intersecar a \( e_2 \) sin usar \( b \), debe contener a \( c \). Luego \( e_3 = \{a,c\} \), y \( e_1, e_2, e_3 \) forman un triángulo sobre \( \{a,b,c\} \). Cualquier otra arista de \( \mathcal{F} \) con un extremo fuera de \( \{a,b,c\} \) dejaría de intersecar a alguna de las tres, por lo que en este caso \( |\mathcal{F}| \leq 3 \).
    \end{itemize}

    En ambos casos \( |\mathcal{F}| \leq 4 \), de modo que no puede haber un conjunto independiente de \( 5 \) vértices. Concluimos que \( S \) es un conjunto independiente máximo del grafo de Petersen.
\end{proof}

\vspace{0.5cm}

\begin{definition}
    Un \textbf{subgrafo} de un grafo \( G \) es un grafo \( H \) tal que \( V(H) \subseteq V(G) \), \( E(H) \subseteq E(G) \), y los extremos de cada arista de \( H \) son los mismos que tiene en \( G \).
\end{definition}

\begin{ejercicioEntregar}{2}
    En cada uno de los grafos de abajo, encontrar un subgrafo bipartito con cantidad máxima de aristas. Justificar por qué no hay un subgrafo bipartito con mayor cantidad de aristas que él y determinar si hay algún otro subgrafo bipartito con la misma cantidad de aristas.
\end{ejercicioEntregar}

\begin{center}
    \begin{tikzpicture}[every node/.style={circle, draw, fill=black, inner sep=1.5pt}]
        \node (l) at (0,0) {};
        \node (d1) at (1,0) {};
        \node (d2) at (2,0) {};
        \node (d3) at (3,0) {};
        \node (d4) at (4,0) {};
        \node (r) at (5,0) {};
        \node (t) at (2,1.2) {};
        \node (b) at (2,-1.2) {};

        \draw[thick] (l) -- (d1) -- (d2) -- (d3) -- (d4) -- (r);
        \draw[thick] (l) -- (t) -- (r);
        \draw[thick] (l) -- (b) -- (r);
        \draw[thick] (t) -- (d2) -- (b);
    \end{tikzpicture}
    \hspace{1.5cm}
    \begin{tikzpicture}[every node/.style={circle, draw, fill=black, inner sep=1.5pt}]
        \node (t1) at (0,1) {};
        \node (t2) at (1,1) {};
        \node (t3) at (2,1) {};
        \node (t4) at (3,1) {};
        \node (b1) at (0,0) {};
        \node (b2) at (1,0) {};
        \node (b3) at (2,0) {};
        \node (b4) at (3,0) {};

        \draw[thick] (t1) -- (t2);
        \draw[thick] (t3) -- (t4);
        \draw[thick] (b1) -- (b2);
        \draw[thick] (b3) -- (b4);
        \draw[thick] (t1) -- (b1);
        \draw[thick] (t2) -- (b2);
        \draw[thick] (t3) -- (b3);
        \draw[thick] (t4) -- (b4);
        \draw[thick] (t2) -- (b3);
        \draw[thick] (t3) -- (b2);
        \draw[thick] (t1) to[bend left=40] (t4);
        \draw[thick] (b1) to[bend right=40] (b4);
    \end{tikzpicture}
\end{center}

\begin{proof}
    Por el \textbf{Teorema 1.2.18} del West \( G \) es bipartito \( \iff \) no contiene ciclos de longitud impar.

    El primer grafo de 11 aristas contiene ciclos de longitud impar (e.g \(l-t-r-d_4-d_3-d_2-d_1-l\)). Eliminemos la mínima cantidad de aristas tal que se rompan todos los ciclos impares. Removiendo \( \{d_4, r\} \) obtenemos un subgrafo bipartito de 10 aristas. Como el grafo original no es bipartito, el subgrafo resultante es maximal. No es único, ya que eliminar \( \{d_3, d_4\} \) o \( \{d_2, d_3\} \) produce subgrafos bipartitos de 10 aristas.

    El segundo grafo de 12 aristas también posee múltiples ciclos impares disjuntos en aristas, lo que significa que eliminar una sola arista no será suficiente para volverlo bipartito. Debemos eliminar al menos dos. Si eliminamos las dos aristas cruzadas del centro (\( \{t_2, b_3\} \) y \( \{t_3, b_2\} \)), el grafo resultante no tiene ciclos impares y por lo tanto es bipartito, conservando 10 aristas. Este subgrafo es maximal. Tampoco es único: eliminando las dos aristas \( \{t_1, t_4\} \) y \( \{b_1, b_4\} \) obtendríamos otro subgrafo bipartito diferente de 10 aristas.

    \begin{center}
        \begin{tikzpicture}[
                vertex/.style={circle, draw, fill=black, inner sep=1.5pt},
                every label/.style={font=\footnotesize}
            ]
            \node[vertex, label=left:\(l\)] (l) at (0,0) {};
            \node[vertex, label=below:\(d_1\)] (d1) at (1,0) {};
            \node[vertex, label=45:\(d_2\)] (d2) at (2,0) {};
            \node[vertex, label=below:\(d_3\)] (d3) at (3,0) {};
            \node[vertex, label=above left:\(d_4\)] (d4) at (4,0) {};
            \node[vertex, label=right:\(r\)] (r) at (5,0) {};
            \node[vertex, label=above:\(t\)] (t) at (2,1.2) {};
            \node[vertex, label=below:\(b\)] (b) at (2,-1.2) {};

            % Se mantiene d4, pero se elimina la arista (d4)--(r)
            \draw[thick] (l) -- (d1) -- (d2) -- (d3) -- (d4);
            \draw[thick] (l) -- (t) -- (r);
            \draw[thick] (l) -- (b) -- (r);
            \draw[thick] (t) -- (d2) -- (b);
        \end{tikzpicture}
        \hspace{1.5cm}
        \begin{tikzpicture}[
                vertex/.style={circle, draw, fill=black, inner sep=1.5pt},
                every label/.style={font=\footnotesize}
            ]
            \node[vertex, label=135:\(t_1\)] (t1) at (0,1) {};
            \node[vertex, label=above:\(t_2\)] (t2) at (1,1) {};
            \node[vertex, label=above:\(t_3\)] (t3) at (2,1) {};
            \node[vertex, label=45:\(t_4\)] (t4) at (3,1) {};
            \node[vertex, label=225:\(b_1\)] (b1) at (0,0) {};
            \node[vertex, label=below:\(b_2\)] (b2) at (1,0) {};
            \node[vertex, label=below:\(b_3\)] (b3) at (2,0) {};
            \node[vertex, label=315:\(b_4\)] (b4) at (3,0) {};

            \draw[thick] (t1) -- (t2);
            \draw[thick] (t3) -- (t4);
            \draw[thick] (b1) -- (b2);
            \draw[thick] (b3) -- (b4);
            \draw[thick] (t1) -- (b1);
            \draw[thick] (t2) -- (b2);
            \draw[thick] (t3) -- (b3);
            \draw[thick] (t4) -- (b4);
            % Aristas cruzadas eliminadas
            \draw[thick] (t1) to[bend left=40] (t4);
            \draw[thick] (b1) to[bend right=40] (b4);
        \end{tikzpicture}
    \end{center}
\end{proof}

Recordemos las nociones de camino (Definición~\ref{def:camino}), ciclo (Definición~\ref{def:ciclo}) y grafo completo (Definición~\ref{def:k4}), todas de la Práctica~II.

\begin{ejercicio}{3}
    Sea \( A \) la familia de los grafos que son caminos, \( B \) la familia de los ciclos, \( C \) la familia de los grafos completos y \( D \) la familia de los grafos bipartitos. Hallar todas las clases de grafos isomorfos que están en la intersección de cada par de esas familias.
\end{ejercicio}
\begin{proof}
    Analicemos cada intersección de familias de grafos:
    \begin{itemize}
        \item \(\mathbf{A \cap B}\): La intersección de caminos y ciclos es vacía. Por definición, en un camino no se repiten vértices, mientras que en un ciclo el primer y último vértice deben coincidir. Notar que un ciclo exige al menos una arista, descartando al grafo trivial.

        \item \(\mathbf{A \cap C}\): La intersección de caminos y grafos completos contiene únicamente a \(K_1\) y \(K_2\). El grafo de un vértice sin aristas (\(K_1\)) es trivialmente un camino y completo. El grafo de dos vértices unidos por una arista (\(K_2\)) también lo es. Cualquier grafo completo \(K_n\) con \(n \geq 3\) tiene vértices de grado \(n-1 \geq 2\) y contiene ciclos (como \(K_3\)), por lo tanto no puede ser un camino.

        \item \(\mathbf{A \cap D}\): La intersección de caminos y grafos bipartitos es la clase de todos los caminos (\(A\)). Como ya dijimos en el ejercicio anterior, un grafo es bipartito si y sólo si no contiene ciclos de longitud impar. Como los caminos no contienen ningún ciclo, todos los caminos son bipartitos.

        \item \(\mathbf{B \cap C}\): La intersección de ciclos y grafos completos contiene únicamente a \(K_3\) (el triángulo). En un ciclo, todo vértice tiene grado exactamente 2. En un grafo completo \(K_n\), todo vértice tiene grado \(n-1\). Igualando ambas condiciones (\(n-1 = 2\)) obtenemos \(n = 3\). Para \(n > 3\), aparecerían diagonales (aristas que no pertenecen al ciclo de la frontera) rompiendo la estructura de ciclo.

        \item \(\mathbf{B \cap D}\): La intersección de ciclos y grafos bipartitos son todos los ciclos de longitud par (\(C_{2k}\) con \(k \geq 2\)). Esto se deduce directamente de la caracterización de los grafos bipartitos (un grafo es bipartito si y sólo si no contiene ciclos de longitud impar).

        \item \(\mathbf{C \cap D}\): La intersección de grafos completos y bipartitos contiene únicamente a \(K_1\) y \(K_2\). Sabemos que un grafo bipartito no puede contener ciclos impares. Cualquier grafo completo \(K_n\) con \(n \geq 3\) contiene a \(K_3\) como subgrafo, que es un ciclo de longitud 3 \( \therefore \) sólo \(K_1\) y \(K_2\) pertenecen a esta intersección.
    \end{itemize}
\end{proof}

\begin{ejercicioEntregar}{4}
    Sea \( G \) el grafo cuyos vértices son las \( k \)-uplas de ceros y unos, con \( x \) adyacente a \( y \) si y sólo si difieren en una posición. Determinar si \( G \) es bipartito.
\end{ejercicioEntregar}
\begin{proof}
    Sea \( k \in \N \). El grafo \(G\) es efectivamente bipartito. Definimos los conjuntos:
    \[
        \begin{aligned}
            X & = \{ v \in V(G) : \text{la cantidad de unos en } v \text{ es par} \},   \\
            Y & = \{ v \in V(G) : \text{la cantidad de unos en } v \text{ es impar} \}.
        \end{aligned}
    \]
    Claramente, \( X \cap Y = \varnothing \) y \(X \cup Y = V(G)\), por lo que conforman una partición válida de los vértices.

    Ahora debemos verificar que no existen aristas con ambos extremos en el mismo conjunto. Supongamos que dos vértices \(u\) y \(v\) son adyacentes en \(G\). Por definición, sus \(k\)-uplas difieren en exactamente una posición.

    Esto implica que la cantidad total de unos en \(u\) y \(v\) difiere exactamente en 1. En consecuencia, si \(u\) tiene una cantidad par de unos, \(v\) tendrá necesariamente una cantidad impar, y viceversa. Es decir, toda arista conecta un vértice de \(X\) con un vértice de \(Y\).

    Por lo tanto, no hay aristas internas en los conjuntos, \((X, Y)\) es una bipartición válida \( \therefore G \) es bipartito.
\end{proof}

\vspace{0.5cm}

Recordemos el algoritmo Bubble Sort (Algoritmo~\ref{alg:bubblesort} de la Práctica~II).

\begin{ejercicio}{5}
    Sea \( G_n \) el grafo cuyos vértices son las permutaciones de \( \{1, \ldots, n\} \), siendo dos permutaciones \( a_1, \ldots, a_n \) y \( b_1, \ldots, b_n \) adyacentes si una puede obtenerse de la otra intercambiando la posición de dos elementos. Probar que \( G_n \) es bipartito. (Pista: para cada permutación \( a \), contar los pares \( i, j \) tales que \( i < j \) y \( a_i > a_j \).)
\end{ejercicio}
\begin{proof}
    Dada una permutación \(a = a_1, \ldots, a_n\), definimos su número de inversiones \( I(a) \) como la cantidad de pares \((i, j)\) tales que \(i < j\) y \(a_i > a_j\).

    El algoritmo Bubble Sort ordena un arreglo realizando únicamente intercambios de elementos \emph{adyacentes}. Es un hecho conocido que cada intercambio de dos elementos adyacentes altera el número de inversiones \(I(a)\) en exactamente \(+1\) o \(-1\). En particular, cada intercambio adyacente cambia la paridad de \(I(a)\).

    En nuestro grafo \(G_n\), dos permutaciones son adyacentes si difieren por un intercambio de \emph{cualesquiera} dos posiciones \(i\) y \(j\) (asumamos \(i < j\)). Un intercambio general entre las posiciones \(i\) y \(j\) se puede simular algorítmicamente mediante una secuencia de intercambios adyacentes (``burbujeando'' el elemento de la posición \(i\) hasta la \(j\), y luego llevando el elemento original de la \(j\) hasta la posición \(i\)). Esto requiere exactamente \((j-i) + (j-i-1) = 2(j-i) - 1\) intercambios adyacentes.

    Como \(2(j-i) - 1\) es siempre un número \textbf{impar}, cualquier arista en \(G_n\) (un intercambio general) equivale a aplicar una cantidad impar de intercambios adyacentes. Como cada uno invierte la paridad de \(I(a)\), deducimos que toda arista de \(G_n\) conecta necesariamente una permutación con un número par de inversiones con una permutación con un número impar de inversiones.

    Si ejecutáramos el \textbf{Algoritmo \ref{alg:bfs-bipartito} (BFS con signos)} (es fácil ver que \( G_n \) es conexo) partiendo del vértice correspondiente a la permutación identidad \( id \) (la cual tiene \(I(id) = 0\), par), el algoritmo le asignaría el número 1. En cada iteración, al cruzar una arista, la paridad de las inversiones necesariamente cambia. Por lo tanto, el BFS asignará sistemáticamente:
    \begin{itemize}
        \item El signo 1 a las permutaciones con \( I(a)\) \textbf{par} (el conjunto \(X\)).
        \item El signo -1 a las permutaciones con \(I(a)\) \textbf{impar} (el conjunto \(Y\)).
    \end{itemize}

    La condición de fallo del algoritmo (línea 9) ocurre si encuentra una arista entre dos vértices con el mismo signo. Pero acabamos de demostrar que toda arista conecta una permutación par con una impar. Por lo tanto, el chequeo de la línea 9 nunca fallará.

    El algoritmo particionará exitosamente el conjunto de vértices de \(G_n\) en dos conjuntos disjuntos \(X\) e \(Y\) sin aristas internas, demostrando que \(G_n\) es bipartito.
\end{proof}

\clearpage

\begin{ejercicio}{6}
    Probar que un grafo \( G \) es bipartito si y sólo si todo subgrafo \( H \) de \( G \) posee un conjunto independiente con al menos la mitad de los vértices que \( H \) tiene.
\end{ejercicio}
\begin{proof}
    (\( \Longrightarrow \)) Sea \(G\) un grafo bipartito con bipartición \((X, Y)\). Para cualquier subgrafo \(H\) de \(G\), los conjuntos \(A = X \cap V(H)\) y \(B = Y \cap V(H)\) forman una partición de \(V(H)\). Como \(X\) e \(Y\) son conjuntos independientes en \(G\), \(A\) y \(B\) también lo son en \(H\). Como \(|A| + |B| = |V(H)|\), al menos uno de los dos conjuntos debe tener un tamaño mayor o igual a \(|V(H)|/2\). Por lo tanto, \(H\) posee un conjunto independiente con al menos la mitad de sus vértices.

    (\( \Longleftarrow \)) Supongamos que todo subgrafo \(H\) de \(G\) posee un conjunto independiente con al menos la mitad de sus vértices, pero que \(G\) no es bipartito.

    Como \(G\) no es bipartito contiene al menos un ciclo de longitud impar. Sea \( H = C_{2k+1} \) dicho subgrafo.

    Analicemos el tamaño máximo de un conjunto independiente en \(H\). Como los vértices forman un ciclo, no podemos elegir dos vértices adyacentes. Lo máximo que podemos hacer es seleccionar vértices alternados. Al ser una cantidad impar de vértices, el conjunto independiente más grande posible tendrá exactamente \(k\) vértices.

    Sin embargo, por nuestra hipótesis inicial, \(H\) debería tener un conjunto independiente de tamaño al menos:
    \[
        \frac{|V(H)|}{2} = \frac{2k+1}{2} = k + \frac{1}{2}
    \]
    Como la cardinalidad de un conjunto debe ser un número entero, el conjunto independiente debería tener al menos \(k+1\) vértices, lo cual es imposible en \(C_{2k+1}\).

    Por lo tanto, \(G\) no puede contener ciclos impares y es bipartito.
\end{proof}

\vspace{0.5cm}

\begin{definition}
    Un \textbf{lazo} (o \emph{loop}) de un grafo \( G \) es una arista cuyos dos extremos coinciden en un mismo vértice.
\end{definition}

\begin{ejercicio}{7}
    Sea \( G \) un grafo sin loops y con \( m \) aristas. Probar que \( G \) posee un subgrafo bipartito con al menos \( \frac{m}{2} \) aristas.
\end{ejercicio}
\begin{proof}
    La idea es la siguiente: de entre todas las particiones posibles de \(V(G)\) en dos conjuntos, elegimos una que maximice la cantidad de aristas ``cruzadas'' (con un extremo en cada conjunto). Esta partición, por ser óptima, ya debe tener al menos \(m/2\) aristas cruzadas: si tuviera menos, podríamos mejorarla moviendo un único vértice de un conjunto al otro, contradiciendo que ya era la mejor posible.

    Formalicemos esta idea. Como \(G\) tiene finitos vértices, hay finitas particiones \((X, Y)\) de \(V(G)\) (permitiendo conjuntos vacíos), así que existe una que maximiza la cantidad de aristas con un extremo en \(X\) y el otro en \(Y\). Fijemos una tal partición \((X, Y)\) y sea \(H\) el subgrafo formado exactamente por esas aristas cruzadas. Por construcción, ninguna arista de \(H\) tiene ambos extremos del mismo lado, así que \(H\) es bipartito con bipartición \((X, Y)\).

    Veamos que todo vértice \(v \in V(G)\) tiene, dentro de \(H\), al menos la mitad de sus aristas originales. Supongamos que no: digamos \(v \in X\) (el caso \(v \in Y\) es análogo) y que \(v\) tiene más vecinos de su mismo lado (\(d_X(v)\) de ellos en \(X\)) que del lado opuesto (\(d_Y(v)\) en \(Y\)), es decir \(d_X(v) > d_Y(v)\). Si moviéramos a \(v\) de \(X\) a \(Y\):
    \begin{itemize}
        \item perderíamos las \(d_Y(v)\) aristas cruzadas que \(v\) ya aportaba (sus vecinos en \(Y\) dejarían de estar del lado opuesto),
        \item pero ganaríamos \(d_X(v)\) aristas cruzadas nuevas (sus vecinos que quedaron en \(X\) ahora sí quedan del lado opuesto a \(v\)).
    \end{itemize}
    Como \(d_X(v) > d_Y(v)\), este cambio aumentaría estrictamente la cantidad de aristas cruzadas, contradiciendo que \((X,Y)\) ya era óptima. Luego, para todo vértice \(v\) debe valer que su cantidad de vecinos del lado opuesto es al menos su cantidad de vecinos del mismo lado, es decir, \(v\) tiene al menos \(\deg(v)/2\) aristas en \(H\):
    \[
        \deg_H(v) \geq \frac{\deg(v)}{2} \qquad \text{para todo } v \in V(G).
    \]

    Sumando esta desigualdad sobre todos los vértices y usando que \(\sum_{v} \deg(v) = 2m\) (cada arista de \(G\) se cuenta dos veces, una por cada extremo):
    \[
        \sum_{v \in V(G)} \deg_H(v) \geq \sum_{v \in V(G)} \frac{\deg(v)}{2} = \frac{1}{2} \sum_{v \in V(G)} \deg(v) = m.
    \]
    Pero el lado izquierdo es, por la misma razón, exactamente \(2|E(H)|\) (cada arista de \(H\) se cuenta dos veces). Se sigue que \(2|E(H)| \geq m\), es decir:
    \[
        |E(H)| \geq \frac{m}{2}.
    \]
    Concluimos que \(H\) es el subgrafo bipartito buscado.
\end{proof}

\vspace{0.5cm}

\begin{definition}
    Definimos la diferencia simétrica de dos conjuntos \(A\) y \(B\) como
    \[
        A \triangle B = (A \setminus B) \cup (B \setminus A).
    \]
    Son los elementos que pertenecen a \( A \) o a \( B \), pero no a ambos.
\end{definition}

\begin{ejercicioEntregar}{8}
    Probar que un grafo bipartito tiene una única bipartición si y sólo si es conexo. Para ello, si se tiene una bipartición \( A, B \), considerar a \( B, A \) como la misma bipartición.
\end{ejercicioEntregar}
\begin{proof}
    (\( \Longrightarrow \)) Procederemos por el contrarrecíproco. Supongamos que \(G\) no es conexo. Entonces, existe una componente conexa de \(G\) cuyo conjunto de vértices llamaremos \(C\), tal que \( C \neq \varnothing \) y \(C \neq V(G)\).

    Sea \((A, B)\) una bipartición válida de \(G\). Consideremos los conjuntos \(A' = A \triangle C\) y \(B' = B \triangle C\). Veamos que \((A', B')\) es también una bipartición de \(G\). Es fácil ver que \( A' \cup B' = V(G) \) y \( A' \cap B' = \varnothing \).

    Toda arista de \(G\) tiene o bien ambos extremos en \(C\), o bien ninguno en \(C\) (pues no hay aristas entre distintas componentes conexas).
    \begin{itemize}
        \item Si una arista no toca \(C\), sus extremos mantienen los mismos conjuntos que en la partición original, por lo que sigue cruzando entre \(A'\) y \(B'\).
        \item Si una arista tiene sus extremos en \(C\), originalmente uno estaba en \(A \cap C\) y el otro en \(B \cap C\). Por definición de la diferencia simétrica, los vértices de \(A \cap C\) pasan a \(B'\) y los de \(B \cap C\) pasan a \(A'\). La arista sigue teniendo sus extremos en conjuntos distintos.
    \end{itemize}
    Por lo tanto, \((A', B')\) es una bipartición. Además, como \( C \neq \varnothing \) y \(C \neq V(G)\), hemos invertido los conjuntos solo para una porción propia del grafo. Esto garantiza que \((A', B')\) no es igual a \((A, B)\) ni a \((B, A)\), contradiciendo así la unicidad. Deducimos entonces que si la bipartición es única, \(G\) debe ser conexo.

    \vspace{0.3cm}

    (\( \Longleftarrow \)) Supongamos que \(G\) es conexo y sean \((A, B)\) y \((A', B')\) dos biparticiones de \(G\). Tomemos un vértice arbitrario \(v_0 \in V(G)\). Sin pérdida de generalidad (intercambiando los nombres de los conjuntos \(A'\) y \(B'\) si fuera necesario), podemos asumir que \(v_0 \in A\) y \(v_0 \in A'\).

    Dado cualquier otro vértice \(u \in V(G)\), como \(G\) es conexo, existe un camino desde \(v_0\) hasta \(u\). Demostraremos por inducción sobre la longitud \(k\) del camino más corto (la distancia) que el conjunto al que pertenece \(u\) está unívocamente determinado.
    \begin{itemize}
        \item \textbf{Caso base (\(k=0\)):} \(u = v_0\), y por nuestra suposición, pertenece a \(A\) y a \(A'\).
        \item \textbf{Paso inductivo:} Supongamos que todo vértice a distancia \(k\) pertenece al mismo conjunto en ambas biparticiones. Sea \(u\) un vértice a distancia \(k+1\), y sea \(w\) su predecesor en el camino mínimo (que está a distancia \(k\)).
              Por hipótesis inductiva, \(w\) tiene la misma asignación en ambas particiones (supongamos que está en \(A\) y \(A'\)). Como \(u\) es adyacente a \(w\) y estamos en un grafo bipartito, \(u\) debe pertenecer forzosamente al conjunto opuesto a \(w\) en ambas particiones. Es decir, \(u \in B\) y \(u \in B'\).
    \end{itemize}
    Como esto vale para todo vértice \(u \in V(G)\), concluimos que \(A = A'\) y \(B = B'\). Por lo tanto, la bipartición es única (salvo el orden de los conjuntos).
\end{proof}

\vspace{0.5cm}

\begin{definition}
    Dado \( n \geq 1 \), notamos \( P_n \) al camino con \( n \) vértices (y, por lo tanto, con \( n-1 \) aristas).
\end{definition}

\begin{definition}
    Un grafo bipartito completo es un grafo bipartito cuya partición \((A, B)\) cumple que todo vértice de \(A\) está conectado con todos los vértices de \(B\) y viceversa.
\end{definition}

\begin{ejercicioEntregar}{9}
    Sea \( G \) un grafo simple y conexo que no tiene a \( P_4 \) ni a \( K_3 \) como subgrafos inducidos. Probar que \( G \) es un grafo bipartito completo.
\end{ejercicioEntregar}
\begin{proof}
    Primero demostraremos que \(G\) es bipartito. Supongamos por el absurdo que \(G\) no lo es, por lo que debe contener al menos un ciclo de longitud impar. Como \(G\) no contiene a \(K_3\) (un ciclo de longitud 3) como subgrafo inducido, el ciclo impar más corto de \(G\) debe tener longitud al menos 5.

    Sea \(C = v_1 v_2 \ldots v_{2k+1} v_1\), con \(k \geq 2\), un ciclo impar de longitud mínima en \(G\). Consideremos los primeros cuatro vértices de este ciclo: \(v_1, v_2, v_3, v_4\). Como no hay un \(P_4\) inducido en \(G\), el camino \(v_1 v_2 v_3 v_4\) no puede ser inducido, por lo que debe existir al menos una arista entre estos vértices. Las únicas posibles son \(v_1v_3\), \(v_2v_4\) o \(v_1v_4\).
    \begin{itemize}
        \item Si existe \(v_1v_3\), los vértices \(v_1, v_2, v_3\) forman un ciclo de longitud 3, es decir, un \(K_3\) inducido, lo cual es una contradicción.
        \item Si existe \(v_2v_4\), análogamente se forma un \(K_3\) inducido en \(v_2, v_3, v_4\), otra contradicción.
        \item Si existe \(v_1v_4\), entonces \(v_1 v_4 v_5 \ldots v_{2k+1} v_1\) forma un ciclo de longitud \((2k+1) - 2 = 2k-1\). Al ser \(k \geq 2\), este es un ciclo impar estrictamente menor que \(C\), contradiciendo la minimalidad de \(C\).
    \end{itemize}
    En cualquier caso llegamos a un absurdo. Por lo tanto, \(G\) no tiene ciclos impares y es un grafo bipartito con alguna partición \((A, B)\).

    Ahora probaremos que \(G\) es bipartito completo. Supongamos por el absurdo que no lo es. Entonces existen vértices \(a \in A\) y \(b \in B\) tales que no son adyacentes (\(ab \notin E(G)\)). Como \(G\) es conexo, existe un camino entre \(a\) y \(b\). Consideremos el \emph{camino mínimo} entre ellos.

    Al estar \(a\) y \(b\) en conjuntos opuestos de la bipartición, cualquier camino entre ellos debe alternar entre \(A\) y \(B\), por lo que su longitud (cantidad de aristas) debe ser impar. Como no son adyacentes, la longitud de este camino mínimo debe ser al menos 3, lo que implica que recorre al menos 4 vértices.

    Sean \(x_1, x_2, x_3, x_4\) los primeros cuatro vértices de este camino mínimo (con \(x_1 = a\)). Por definición de camino mínimo, no puede haber aristas que conecten vértices no consecutivos dentro de él (de lo contrario, habría un camino más corto). Por lo tanto, el subgrafo inducido por \( \{x_1, x_2, x_3, x_4\} \) es exactamente un \( P_4 \), absurdo.

    Por lo tanto, \(G\) es un grafo bipartito completo.
\end{proof}

\vspace{0.5cm}

\begin{definition}
    Dados grafos \( G_1, \ldots, G_k \), la \textbf{unión} \( G_1 \cup \cdots \cup G_k \) es el grafo con conjunto de vértices \( V(G_1) \cup \cdots \cup V(G_k) \) y conjunto de aristas \( E(G_1) \cup \cdots \cup E(G_k) \).
\end{definition}

\begin{ejercicioEntregar}{10}
    \begin{enumerate}
        \item Dados \( n \) y \( k \) números naturales tales que \( n \leq 2^k \), codificar los vértices de \( K_n \) con \( k \)-uplas y usar esto para construir \( k \) grafos bipartitos cuya unión es \( K_n \).
        \item Recíprocamente, si \( K_n \) es la unión de \( k \) grafos bipartitos \( G_1, \ldots, G_k \), codificar convenientemente los vértices de \( K_n \) con \( k \)-uplas binarias y probar así que \( n \leq 2^k \).
    \end{enumerate}
\end{ejercicioEntregar}
\begin{proof}
    (\( \Longrightarrow \)) Como \(n \leq 2^k\), podemos asignar a cada vértice de \(K_n\) una \(k\)-upla binaria \((x_1, \ldots, x_k) \in {\{0,1\}}^k\) única.
    Para cada \( i \in \{1, \ldots, k\} \), definimos un grafo \(G_i\) sobre el mismo conjunto de vértices \(V(K_n)\). Particionamos los vértices en dos conjuntos según el valor de su \(i\)-ésima coordenada:
    \begin{align*}
        A_i & = \{ v \in V(K_n) : \text{la } i\text{-ésima coordenada de } v \text{ es } 0 \}, \\
        B_i & = \{ v \in V(K_n) : \text{la } i\text{-ésima coordenada de } v \text{ es } 1 \}
    \end{align*}
    Declaramos que \(G_i\) contiene exactamente todas las aristas de \(K_n\) que conectan un vértice de \(A_i\) con uno de \(B_i\). Por construcción, no hay aristas internas en los conjuntos, luego cada \(G_i\) es bipartito.

    Para ver que \(K_n = \bigcup_{i=1}^k G_i\), tomemos dos vértices distintos cualesquiera \(u, v \in V(K_n)\). Al tener \(k\)-uplas únicas, deben diferir en al menos una coordenada, digamos la coordenada \(j\). Esto implica que uno pertenece a \(A_j\) y el otro a \(B_j\). Luego, la arista \(uv\) pertenece a \(G_j\). Como toda arista de \(K_n\) está presente en al menos un \(G_i\), la unión de todos los subgrafos cubre \(K_n\).

    \vspace{0.3cm}

    (\( \Longleftarrow \)) Supongamos que \(K_n = \bigcup_{i=1}^k G_i\), donde cada \(G_i\) es un grafo bipartito.
    Cada \(G_i\) admite una bipartición. Si hay vértices de \(K_n\) que no tienen aristas en un determinado \(G_i\) (vértices aislados en ese subgrafo), podemos ubicarlos arbitrariamente en cualquiera de los dos conjuntos de su bipartición sin generar aristas internas. De este modo, podemos garantizar que cada \(G_i\) posee una bipartición \((A_i, B_i)\) tal que \(A_i \cup B_i = V(K_n)\).

    Asignamos a cada vértice \(v \in V(K_n)\) una \(k\)-upla binaria \((x_1, \ldots, x_k)\), usando la siguiente regla para cada coordenada \(i\):
    \[
        x_i = 0 \text{ si } v \in A_i, \quad x_i = 1 \text{ si } v \in B_i
    \]
    Supongamos por reducción al absurdo que \(n > 2^k\). Por el principio del palomar, al haber estrictamente más vértices que \(k\)-uplas binarias posibles (que son exactamente \(2^k\)), al menos dos vértices distintos \(u\) y \(v\) deben recibir la misma tupla.

    Tener el mismo código significa que, para todo \(i = 1, \ldots, k\), ambos vértices caen del mismo lado de la partición (ambos en \(A_i\) o ambos en \(B_i\)). Como los conjuntos son independientes, la arista \(uv\) no pertenece a \(G_i\) para ningún \(i\). En consecuencia, la arista \(uv\) no pertenece a la unión \(\bigcup_{i=1}^k G_i\).

    Pero esto es una contradicción, ya que la unión de los subgrafos compone el grafo completo \(K_n\), donde la arista \(uv\) obligatoriamente debe existir.

    Esta contradicción proviene de suponer \(n > 2^k\). Concluimos entonces que todas las tuplas de los vértices deben ser distintas, por lo que \(n \leq 2^k\).
\end{proof}

\vspace{0.5cm}

\begin{ejercicio}{11}
    Dado un grafo \( G \) conexo y un vértice inicial \( v \), consideremos la siguiente variante de BFS donde a los vértices les asignamos un 1 o un -1.
\end{ejercicio}

\begin{algorithm}[h]
    \caption{BFS con signos}
    \label{alg:bfs-bipartito}
    \begin{algorithmic}[1]
        \Procedure{BFSBipartito}{$G, v$}
        \State \textbf{Entrada:} Un grafo conexo $G$ y un vértice inicial $v$.
        \State \textbf{Salida:} Una asignación de signos $\pm 1$ a los vértices de $G$, o la conclusión de que $G$ no es bipartito.
        \vspace{0.2cm}
        \State Introducir a $v$ en la cola y asignarle el número 1.
        \While{la cola no esté vacía}
        \State Considerar el primer vértice $w$ de la cola y removerlo de ella.
        \For{cada vecino $x$ de $w$}
        \If{$x$ ya tiene un número asignado}
        \State Chequear que sea el opuesto al de $w$. En caso contrario, finalizar y concluir que $G$ no es bipartito.
        \Else
        \State Asignarle a $x$ el número opuesto al de $w$ y colocar a $x$ en la cola.
        \EndIf
        \EndFor
        \EndWhile
        \EndProcedure
    \end{algorithmic}
\end{algorithm}

Demostrar que este algoritmo funciona, es decir, si \( G \) es bipartito hallará una partición del conjunto de vértices que lo demuestra y, si \( G \) no es bipartito, alcanzará tal conclusión.
\begin{proof}
    Sea \(G\) un grafo conexo (y finito). Analizaremos los dos posibles resultados del algoritmo:

    Si el algoritmo termina sin encontrar contradicciones, significa que ha asignado signos a todos (gracias a la conexidad) los vértices de manera consistente. En este caso, podemos definir los conjuntos \( A = \{ u \in V(G) : \text{sign}(u) = 1 \} \) y \( B = \{ u \in V(G) : \text{sign}(u) = -1 \} \), obteniendo así una bipartición válida de \(G\). El propio algoritmo nos dice que todos los vecinos de un vértice tienen signos opuestos, lo que garantiza que no existen aristas dentro de \(A\) ni dentro de \(B\).

    Si el algoritmo falla se concluye que \(G\) no es bipartito. Esto ocurre únicamente si al explorar los vecinos de un vértice \(w\), se encuentra un vecino \(x\) que ya tiene asignado el mismo signo que \(w\).

    Notemos que el signo que el BFS le asigna a un vértice \(u\) depende exclusivamente de su distancia al vértice inicial \(v\). Específicamente, el signo alterna en cada paso, por lo que \(\text{sign}(u) = {(-1)}^{d(v, u)}\).
    Si \(w\) y \(x\) tienen el mismo signo, esto implica que sus distancias al vértice inicial, \(d(v, w)\) y \(d(v, x)\), tienen la misma paridad.

    Consideremos el recorrido cerrado formado por: el camino más corto de \(v\) a \(w\) (longitud \(d(v,w)\)), la arista \(wx\) (longitud 1), y el camino más corto de \(x\) de regreso a \(v\) (longitud \(d(v,x)\)). La longitud total de este recorrido cerrado es:
    \[
        L = d(v, w) + 1 + d(v, x)
    \]
    Como \(d(v, w)\) y \(d(v, x)\) tienen la misma paridad, su suma es un número par. Al sumarle 1 (la arista \(wx\)), obtenemos que \(L\) es impar. Sabemos que un recorrido cerrado de longitud impar garantiza la existencia de un \textbf{ciclo de longitud impar} en \(G \implies G \) no es bipartito.
\end{proof}
\clearpage

\section{Práctica IV}

\begin{definition}\label{def:kmn}
    Dados \( m, n \geq 1 \), el \textbf{grafo bipartito completo} \( K_{m,n} \) es el grafo bipartito con bipartición \( A, B \) tal que \( |A| = m \), \( |B| = n \), y cada vértice de \( A \) es adyacente a cada vértice de \( B \) (en particular, \( |E(K_{m,n})| = mn \)).
\end{definition}

\begin{definition}
    Un \textbf{circuito euleriano} de un grafo \( G \) es una cadena cerrada que contiene a todas las aristas de \( G \), es decir, que recorre cada arista de \( G \) exactamente una vez. Un grafo se dice \textbf{euleriano} si posee un circuito euleriano.
\end{definition}

\begin{ejercicio}{1}
    Determinar los posibles valores de \( m \) y \( n \) para los que el grafo bipartito completo \( K_{m,n} \) sea euleriano.
\end{ejercicio}
\begin{proof}
    Recordemos que un grafo conexo posee un circuito euleriano si y sólo si todos sus vértices tienen grado par.

    Consideremos el grafo bipartito completo y conexo \(K_{m,n}\) con particiones \(A\) y \(B\), donde \(|A| = m\) y \(|B| = n\). Por definición, cada vértice de un conjunto está adyacente a todos los vértices del otro conjunto. Entonces, tenemos que:
    \begin{itemize}
        \item Todo vértice \(v \in A\) está conectado a los \(n\) vértices de \(B\), por lo que \(\text{deg}(v) = n\).
        \item Todo vértice \(u \in B\) está conectado a los \(m\) vértices de \(A\), por lo que \(\text{deg}(u) = m\).
    \end{itemize}

    Para que exista un circuito euleriano, la condición necesaria y suficiente es que todos los vértices del grafo tengan grado par. Según lo anterior, esto equivale a exigir que tanto \(n\) como \(m\) sean números pares.

    Concluimos entonces que \(K_{m,n}\) es euleriano si y sólo si \(m\) y \(n\) son ambos números pares.
\end{proof}

\vspace{0.5cm}

\begin{ejercicio}{2}
    Determinar si las siguientes afirmaciones son verdaderas o falsas:
    \begin{itemize}
        \item Todo grafo bipartito euleriano tiene un número par de aristas.
        \item Todo grafo simple conexo y euleriano con un número par de vértices tiene un número par de aristas.
        \item Si \( G \) es euleriano y \( v \) es un vértice de \( G \), entonces para cualquier par de aristas \( e \) y \( f \) que inciden en \( v \) existe un circuito euleriano en el cual \( e \) y \( f \) aparecen consecutivamente.
    \end{itemize}
\end{ejercicio}
\begin{proof}
    Las proposiciones son: verdadera, falsa y falsa respectivamente.
    \begin{itemize}
        \item Sea \((A, B)\) la bipartición del grafo \(G\). Como \(G\) es euleriano, el grado de todo vértice es par. En un grafo bipartito, toda arista tiene exactamente un extremo en el conjunto \(A\). Por lo tanto, si sumamos los grados de todos los vértices en \(A\), estaremos contando cada arista de \(G\) exactamente una vez. Es decir, la cantidad total de aristas \(|E|\) es:
              \[
                  |E| = \sum_{u \in A} \text{deg}(u)
              \]
              Como cada vértice \(u \in A\) tiene grado par (\(\text{deg}(u) = 2k\) para algún \( k \in \N \)), la suma está compuesta exclusivamente por números pares, por lo que el total de aristas \(|E|\) es un número par.

        \item Como contraejemplo consideremos un grafo formado por la unión de un ciclo de 3 vértices (\(C_3\)) y un ciclo de 4 vértices (\(C_4\)) que comparten exactamente un vértice \(v\) (y ninguna otra arista ni vértice).
              \begin{itemize}
                  \item El grafo es conexo.
                  \item El vértice compartido \(v\) tiene grado 4, y los restantes vértices tienen grado 2. Como todos los grados son pares, el grafo es euleriano.
                  \item La cantidad total de vértices es 6.
                  \item La cantidad total de aristas es  7.
              \end{itemize}
              Esto contradice la afirmación.

        \item Como contraejemplo consideremos el grafo formado por dos ciclos de 3 vértices (\(T_1\) y \(T_2\)) que comparten un único vértice \(v\). Este grafo es conexo y todos sus vértices tienen grado par, por lo que es euleriano.
              Sean \(e\) y \(f\) las dos aristas del ciclo \(T_1\) que inciden en el vértice \(v\).

              Supongamos por el absurdo que existe un circuito euleriano donde \(e\) y \(f\) aparecen consecutivamente.
              Esto significa que el circuito llega a \(v\) por \(e\) e inmediatamente sale por \(f\). Como los otros extremos de \(e\) y \(f\) (llamémoslos \(x\) e \(y\)) tienen grado 2 y están conectados entre sí por la tercera arista de \(T_1\), el circuito, al llegar a \(y\) a través de \(f\), se ve forzado a tomar la arista \(yx\) y luego regresar a \(v\).

              En este punto, el recorrido se queda sin aristas sin visitar en \(T_1\) para poder salir de allí. Como el circuito nunca cruzó hacia \(T_2\), no puede ser un circuito euleriano (no recorrió todas las aristas). Por lo tanto, \(e\) y \(f\) no pueden ser recorridas consecutivamente.

    \end{itemize}
\end{proof}

\vspace{0.5cm}

\begin{ejercicio}{3}
    \begin{enumerate}
        \item Sea \( G \) un grafo con conjunto de aristas no vacío y tal que todos sus vértices son de grado par. Probar que \( G \) posee un ciclo.
        \item Probar que un grafo conexo es euleriano si y sólo si puede expresarse como una unión \( C_1 \cup \cdots \cup C_n \) de ciclos que no comparten aristas.
    \end{enumerate}
\end{ejercicio}
\begin{proof}
    \begin{enumerate}
        \item Como el conjunto de aristas es no vacío, existe al menos una componente conexa de \(G\) con aristas. Al ser todos los grados pares, los vértices de esta componente deben tener grado mayor o igual a 2 (no hay vértices de grado 1).

              Sea \(P = v_0, v_1, \ldots, v_k\) un camino de longitud máxima en esta componente. Consideremos el último vértice del camino, \(v_k\). Como su grado es al menos 2, es incidente a al menos otra arista además de \(v_{k-1}v_k\). Por la maximalidad del camino \(P\), \(v_k\) no puede ser adyacente a ningún vértice fuera de \(P\), ya que en ese caso podríamos extender el camino y \(P\) no sería maximal.

              Por lo tanto, todos los vecinos de \(v_k\) ya pertenecen a \(P\). Sea \(v_i\) uno de esos vecinos con \(i < k-1\). La arista \(v_k v_i\) junto con la porción del camino \(P\) que va desde \(v_i\) hasta \(v_k\) forman un ciclo. Luego, \(G\) posee al menos un ciclo.

        \item (\( \Longrightarrow \)) Supongamos que \(G\) es conexo y euleriano. Procederemos por inducción fuerte sobre la cantidad de aristas \(m\).
              Como \(G\) es euleriano, todos sus vértices tienen grado par. Por el ítem 1, al tener aristas y grados pares, \(G\) posee al menos un ciclo \(C_1\).
              Si \(G = C_1\), la prueba termina. Si no, consideremos el grafo \(G' = G - E(C_1)\) que se obtiene al eliminar las aristas de \(C_1\).

              Al quitar un ciclo, el grado de todo vértice que pertenece al ciclo se reduce exactamente en 2, y el de los demás vértices no se altera. Por lo tanto, todos los vértices de \(G'\) siguen teniendo grado par. Aunque \(G'\) podría ya no ser conexo, cada una de sus componentes conexas tiene vértices de grado par y estrictamente menos de \(m\) aristas. Por hipótesis inductiva, cada componente con aristas se puede descomponer en una unión de ciclos disjuntos en aristas. Juntando todos estos ciclos con nuestro \(C_1\) original, obtenemos una colección de ciclos \(C_1, \ldots, C_k\) disjuntos en aristas cuya unión es \(G\).

              (\( \Longleftarrow \)) Supongamos que el grafo conexo \(G\) puede expresarse como una unión de ciclos disjuntos en aristas \(G = C_1 \cup \cdots \cup C_n\).
              Queremos ver que \(G\) es euleriano. Para ello, basta con calcular el grado de cada vértice. Si tomamos un vértice arbitrario \(v \in V(G)\), su grado total en \(G\) será la suma de los grados que le aporta cada ciclo en el que participa.
              Como todo vértice en un ciclo tiene grado exactamente 2 dentro de ese ciclo, cada vez que \(v\) pertenece a algún \(C_i\), este ciclo contribuye con un \(+2\) al grado total de \(v\).
              Como los ciclos no comparten aristas, la contribución es aditiva, por lo que el grado de \(v\) será \(2 \times (\text{cantidad de ciclos que contienen a } v)\), que es un número par.
              Al ser un grafo conexo donde todos los vértices tienen grado par, por el teorema fundamental, \(G\) es un grafo euleriano.
    \end{enumerate}
\end{proof}

\vspace{0.5cm}

\begin{definition}
    Dos circuitos eulerianos se dicen equivalentes si tienen el mismo orden cíclico
    de sus aristas (para toda arista, sus aristas vecinas son las mismas en ambos circuitos, considerándose en cada circuito a la primera arista vecina de la
    última). Un ciclo, por ejemplo, tiene una única clase de equivalencia de circuitos eulerianos.
\end{definition}

\begin{ejercicio}{4}
    ¿Cuántas clases de equivalencia de circuitos eulerianos hay en el siguiente grafo?
\end{ejercicio}

\begin{center}
    \begin{tikzpicture}[vertice/.style={circle, draw, fill=black, inner sep=1.5pt}]
        \node[vertice, label=above:\(y\)] (t) at (0,1.5) {};
        \node[vertice, label=below:\(x\)] (b) at (0,-1.5) {};
        \node[vertice, label=left:\(a\)] (l) at (-2.2,0) {};
        \node[vertice, label=right:\(d\)] (r) at (2.2,0) {};
        \node[vertice, label=left:\(b\)] (i1) at (-0.7,0) {};
        \node[vertice, label=right:\(c\)] (i2) at (0.7,0) {};

        \draw[thick] (t) -- (l) -- (b) -- (r) -- (t);
        \draw[thick] (t) -- (i1) -- (b);
        \draw[thick] (t) -- (i2) -- (b);
    \end{tikzpicture}
\end{center}
\begin{proof}
    El grafo es conexo y todos sus vértices tienen grado par (\(y\) y \(x\) tienen grado 4; \(a, b, c\) y \(d\) tienen grado 2), por lo que es euleriano.

    Notemos que \(y\) y \(x\) están unidos por exactamente cuatro caminos internamente disjuntos de longitud 2: \(y\text{-}a\text{-}x\), \(y\text{-}b\text{-}x\), \(y\text{-}c\text{-}x\) y \(y\text{-}d\text{-}x\). Como los vértices \(a, b, c, d\) tienen grado 2, cualquier circuito euleriano que entre a uno de ellos debe salir inmediatamente. Por lo tanto, cada camino se comporta estructuralmente como una única arista indivisible entre \(y\) y \(x\).

    Cualquier circuito euleriano asume la forma de una secuencia alternada entre los polos:
    \[
        y \xrightarrow{\,e_{\sigma(1)}\,} x \xrightarrow{\,e_{\sigma(2)}\,} y \xrightarrow{\,e_{\sigma(3)}\,} x \xrightarrow{\,e_{\sigma(4)}\,} y,
    \]
    para alguna permutación \( \sigma \) de \( \{a,b,c,d\} \). Dado que no hay restricciones de adyacencia adicionales, cualquiera de las \(4! = 24\) permutaciones genera un circuito euleriano válido.

    Para agrupar estas 24 secuencias en clases de equivalencia, debemos considerar cuándo dos secuencias representan el mismo recorrido. Dos permutaciones son equivalentes si difieren solo en el vértice de inicio o en el sentido del recorrido.

    Como los caminos alternan estrictamente entre \(y\) y \(x\), una rotación impar de la secuencia (e.g pasar de \((a,b,c,d)\) a \((b,c,d,a)\)) invierte los roles de \(x\) e \(y\), alterando qué caminos se emparejan en cada vértice y generando así un circuito distinto. En consecuencia, cada circuito está representado por exactamente 4 permutaciones equivalentes:
    \begin{itemize}
        \item La secuencia original.
        \item Su rotación de 2 posiciones (que preserva los emparejamientos en \(y\) y \(x\)).
        \item Su reflexión (recorrer el circuito en sentido opuesto).
        \item La rotación de 2 posiciones de su reflexión.
    \end{itemize}

    Dividiendo el total de secuencias posibles por la cantidad de representaciones equivalentes, obtenemos la cantidad de clases de equivalencia:
    \( \frac{24}{4} = 6 \)

    Concluimos que el grafo tiene exactamente 6 clases de equivalencia de circuitos eulerianos.
\end{proof}

\vspace{0.5cm}

\begin{definition}
    Una cadena de un grafo \( G \) es una secuencia de aristas \( (e_1, e_2, \ldots, e_k) \) tal que cada par de aristas consecutivas \( e_i \) y \( e_{i+1} \) comparten un vértice, y ninguna arista se repite en la secuencia.
\end{definition}

\begin{definition}
    Una \textbf{descomposición en cadenas} de un grafo \( G \) es una colección de cadenas \( R_1, \ldots, R_k \) de \( G \) que no comparten aristas entre sí y tales que toda arista de \( G \) pertenece a alguna \( R_i \).
\end{definition}

\begin{ejercicio}{5}
    ¿Cuál es el mínimo número \( k \) de cadenas que se necesitan para descomponer al grafo de Petersen? ¿Puede descomponerse también en \( k \) cadenas que sean caminos?
\end{ejercicio}
\begin{proof}
    Sea \( G \) el grafo de Petersen. Por definición, \(|V(G)| = 10\) y \( \forall v \in V(G)\), \(\text{deg}(v) = 3\).

    Sea \( P = \{P_1, P_2, \dots, P_k\} \) una partición del conjunto de aristas \(E(G)\) en \(k\) cadenas.

    Para cualquier vértice \(v \in V(G)\), su grado total en \(G\) es exactamente la suma de los grados que tiene en cada subgrafo inducido por las cadenas \(P_i\). Es decir:
    \[
        \text{deg}_G(v) = 3 = \sum_{i=1}^k \text{deg}_{P_i}(v)
    \]

    Para cada cadena \(P_i\), consideremos la paridad del grado que aporta a un vértice \(v\):
    \begin{itemize}
        \item Si \(P_i\) es una cadena \emph{cerrada}, entonces \(\text{deg}_{P_i}(v)\) es par para todo \(v\), ya que cada paso por \(v\) como vértice interno aporta \(2\) a su grado.
        \item Si \(P_i\) es una cadena con extremos distintos \(a_i \ne b_i\), entonces \(\text{deg}_{P_i}(v)\) es impar si y solo si \( v \in \{a_i, b_i\} \), pues cada paso interno aporta \(2\) y solo el pasaje por un extremo aporta \(1\).
    \end{itemize}

    Como \(\deg_G(v) = 3\) es impar para todo \(v\), y esta suma es la de los \(\deg_{P_i}(v)\), debe haber una cantidad impar (en particular, al menos una) de sumandos impares. Por lo anterior, esto significa que todo vértice \(v \in V(G)\) debe ser extremo de \emph{al menos una} cadena abierta \(P_i \in P\).

    Esto nos dice que los 10 vértices de \(G\) deben ser el extremo de al menos una cadena en la descomposición de \(G\).

    Sea \( \alpha \) el número total de extremos disponibles en la partición \(P\). Dado que cada cadena tiene a lo sumo 2 extremos, el número total de extremos es:
    \[
        \alpha \le 2k
    \]

    Como necesitamos cubrir los 10 vértices de \(G\) con al menos un extremo cada uno, el número total de extremos en la partición debe ser mayor o igual a 10. Por lo tanto:
    \begin{align*}
        2k & \ge 10 \\
        k  & \ge 5
    \end{align*}
    Queda demostrado que se requieren al menos 5 cadenas para realizar la descomposición.

    Veamos que la cota \(k=5\) se alcanza, incluso con caminos simples, lo que responde también la segunda pregunta.
    \begin{center}
        \begin{tikzpicture}[
                vertex/.style={circle, draw, fill=black, inner sep=1.5pt, outer sep=0pt},
                legend/.style={font=\small},
            ]

            \def\R{2.8}
            \def\r{1.3}

            % --- Vértices ---
            \node[vertex] (o1) at (90:\R) {};
            \node[above=3pt] at (o1) {12};

            \node[vertex] (o2) at (18:\R) {};
            \node[above right=2pt] at (o2) {34};

            \node[vertex] (o3) at (306:\R) {};
            \node[below right=2pt] at (o3) {51};

            \node[vertex] (o4) at (234:\R) {};
            \node[below left=2pt] at (o4) {23};

            \node[vertex] (o5) at (162:\R) {};
            \node[above left=2pt] at (o5) {45};

            \node[vertex] (i1) at (90:\r) {};
            \node[right=3pt] at (i1) {35};

            \node[vertex] (i2) at (18:\r) {};
            \node[below right=2pt] at (i2) {52};

            \node[vertex] (i3) at (306:\r) {};
            \node[left=3pt] at (i3) {24};

            \node[vertex] (i4) at (234:\r) {};
            \node[below right] at (i4) {41};

            \node[vertex] (i5) at (162:\r) {};
            \node[above right=2pt] at (i5) {13};

            % --- Cadenas coloreadas (cubren las 15 aristas exactamente una vez) ---
            \draw[red, very thick]    (o2) -- (o1) -- (i1) -- (i3); % C1: 34-12-35-24
            \draw[blue, very thick]   (o3) -- (o2) -- (i2) -- (i4); % C2: 51-34-52-41
            \draw[teal, very thick]   (o4) -- (o3) -- (i3) -- (i5); % C3: 23-51-24-13
            \draw[orange, very thick] (o5) -- (o4) -- (i4) -- (i1); % C4: 45-23-41-35
            \draw[violet, very thick] (o1) -- (o5) -- (i5) -- (i2); % C5: 12-45-13-52

            % --- Leyenda ---
            \begin{scope}[shift={(4.6,0)}, legend]
                \draw[red, very thick] (0,1.6) -- (0.8,1.6);
                \node[right] at (0.8,1.6) {$C_1: 34\text{-}12\text{-}35\text{-}24$};

                \draw[blue, very thick] (0,0.8) -- (0.8,0.8);
                \node[right] at (0.8,0.8) {$C_2: 51\text{-}34\text{-}52\text{-}41$};

                \draw[teal, very thick] (0,0) -- (0.8,0);
                \node[right] at (0.8,0) {$C_3: 23\text{-}51\text{-}24\text{-}13$};

                \draw[orange, very thick] (0,-0.8) -- (0.8,-0.8);
                \node[right] at (0.8,-0.8) {$C_4: 45\text{-}23\text{-}41\text{-}35$};

                \draw[violet, very thick] (0,-1.6) -- (0.8,-1.6);
                \node[right] at (0.8,-1.6) {$C_5: 12\text{-}45\text{-}13\text{-}52$};
            \end{scope}

        \end{tikzpicture}
    \end{center}

    Lo anterior muestra que el número mínimo de cadenas necesarias para cubrir todas las aristas del grafo exactamente una vez es cinco y que se alcanza con una colección de caminos disjuntos que cubren todas las aristas.
\end{proof}

\vspace{0.5cm}

\begin{definition}
    El \textbf{problema del cartero chino} consiste en, dado un grafo conexo \( G \), hallar un recorrido cerrado de longitud mínima que contenga a todas las aristas de \( G \), pudiendo repetir aristas de ser necesario.
\end{definition}

\begin{definition}
    El \textbf{grafo de los puentes de Königsberg} es el grafo múltiple cuyos vértices son las cuatro regiones de tierra de la ciudad de Königsberg —las riberas norte (\( N \)) y sur (\( S \)) del río, y las islas \( K \) y \( O \)— y cuyas siete aristas se corresponden con los siete puentes que las conectaban: dos puentes entre \( N \) y \( K \), dos puentes entre \( S \) y \( K \), un puente entre \( K \) y \( O \), un puente entre \( N \) y \( O \), y un puente entre \( S \) y \( O \).
\end{definition}

\begin{center}
    \begin{tikzpicture}[every node/.style={circle, draw, fill=black, inner sep=1.5pt}]
        \node[label=above:$N$] (n) at (0,1.5) {};
        \node[label=below:$S$] (s) at (0,-1.5) {};
        \node[label=left:$K$] (k) at (-2,0) {};
        \node[label=right:$O$] (o) at (2,0) {};

        \draw[thick] (n) to[bend left=20] (k);
        \draw[thick] (n) to[bend right=20] (k);
        \draw[thick] (s) to[bend left=20] (k);
        \draw[thick] (s) to[bend right=20] (k);
        \draw[thick] (n) -- (o);
        \draw[thick] (s) -- (o);
        \draw[thick] (k) -- (o);
    \end{tikzpicture}
    \\[.25cm]
    El grafo de los puentes de Königsberg.
\end{center}

\begin{ejercicio}{6}
    Resolver el problema del cartero chino en el grafo de los puentes de Königsberg.
\end{ejercicio}
\begin{proof}
    Notemos que el problema equivale a agregar la mínima cantidad de aristas posibles tal que todos los vértices resulten de grado par. Esto se puede hacer de la siguiente forma:

    \begin{center}
        \begin{tikzpicture}[every node/.style={circle, draw, fill=black, inner sep=1.5pt}]
            \node[label=above:\(N\)] (n) at (0,1.5) {};
            \node[label=below:\(S\)] (s) at (0,-1.5) {};
            \node[label=left:\(K\)] (k) at (-2,0) {};
            \node[label=right:\(O\)] (o) at (2,0) {};

            \draw[thick] (n) to[bend left=20] (k);
            \draw[thick] (n) to[bend right=20] (k);
            \draw[thick] (s) to[bend left=20] (k);
            \draw[thick] (s) to[bend right=20] (k);
            \draw[thick] (n) -- (o);
            \draw[thick] (s) -- (o);
            \draw[thick] (k) -- (o);
            \draw[red, thick] (n) to[bend left=20] (o);
            \draw[red, thick] (s) to[bend left=60] (k);
        \end{tikzpicture}
    \end{center}

    Logrando que: \(\text{deg}(K) = 6\), \(\text{deg}(N) = \text{deg}(S) = \text{deg}(O) = 4\). Como todos los vértices tienen grado par y el grafo es claramente conexo se puede encontrar un circuito euleriano.

    Finalmente el camino buscado es el recorrido euleriano que pasa por las 9 aristas exactamente una vez, dado por:
    \[
        N \to K \to S \to K \to S \to O \to N \to O \to K \to N
    \]
\end{proof}

\clearpage

\section{Práctica V}
\documentclass{article}

\usepackage[margin=1in]{geometry}
\usepackage{amsmath,amsthm,amssymb}
\usepackage[spanish]{babel}
\usepackage{enumitem}
\usepackage{graphicx}
\usepackage{tikz}
\usepackage{algorithm}
\usepackage{algpseudocode}
\usetikzlibrary{calc}

\theoremstyle{plain}
\newtheorem{proposicion}{Proposición}[section]

\theoremstyle{definition}
\newtheorem{definition}{Definición}[section]

% Number sets
\newcommand{\R}{\mathbb{R}}
\newcommand{\Z}{\mathbb{Z}}
\newcommand{\N}{\mathbb{N}}
\newcommand{\Q}{\mathbb{Q}}
\newcommand{\C}{\mathbb{C}}

% Exercise environment
\newenvironment{ejercicio}[1]{\begin{trivlist}
\item[\hskip \labelsep {\bfseries Ejercicio}\hskip \labelsep {\bfseries #1.}]}{\end{trivlist}}

\setlist[enumerate,1]{label=(\alph*), leftmargin=*, itemsep=2pt}

\begin{document}

\title{Teoría de Grafos}
\author{Bustos Jordi\\Práctica V}

\maketitle

\begin{definition}
    Un \textbf{camino} de un grafo \( G \) es una secuencia de vértices \( v_1, v_2, \dots, v_k \) tal que \( v_i \) es adyacente a \( v_{i+1} \) para todo \( i = 1, \dots, k-1 \) y no se repiten vértices.
\end{definition}

\begin{definition}
    Un \textbf{ciclo} de un grafo \( G \) es un camino cerrado que no repite vértices, excepto el vértice inicial y final.
\end{definition}

\begin{definition}
    Un \textbf{ciclo hamiltoniano} de un grafo \( G \) es un ciclo que contiene a todos los vértices de \( G \). Un grafo se dice \textbf{hamiltoniano} si posee un ciclo hamiltoniano.
\end{definition}

\begin{ejercicio}{1}
    Determinar si el siguiente grafo es hamiltoniano.
\end{ejercicio}

\begin{center}
    \begin{tikzpicture}[vertice/.style={circle, draw, fill=black, inner sep=1.5pt}]
        \def\R{2.3}
        \def\r{1.05}
        \foreach \k in {0,...,7} {
                \node[vertice] (o\k) at ({22.5+45*\k}:\R) {};
                \node[vertice] (i\k) at ({22.5+45*\k}:\r) {};
            }
        \foreach \k in {0,...,7} {
                \pgfmathtruncatemacro{\next}{mod(\k+1,8)}
                \draw[thick] (o\k) -- (o\next);
                \draw[thick] (o\k) -- (i\k);
            }
        \foreach \k in {0,...,7} {
                \pgfmathtruncatemacro{\next}{mod(\k+2,8)}
                \draw[thick] (i\k) -- (i\next);
            }
    \end{tikzpicture}
\end{center}

\begin{proof}
    Etiquetamos \( u_0, \ldots, u_7 \) los vértices exteriores en sentido antihorario y \( v_0, \ldots, v_7 \) los interiores también en sentido antihorario.
    El camino:
    \[ u_0 - v_0 - v_2 - u_2 - u_1 - v_1 - v_3 - u_3 - u_4 - v_4 - v_6 - u_6 - u_5 - v_5 - v_7 - u_7 - u_0 \]
    Es un ciclo hamiltoniano.
\end{proof}

\vspace{0.5cm}

\begin{ejercicio}{2}
    Hallar un grafo simple y conexo, con a lo sumo 6 vértices y al menos una arista, que verifique cada una de las siguientes condiciones:
    \begin{enumerate}
        \item \( G \) es euleriano y hamiltoniano.
        \item \( G \) es euleriano, pero no hamiltoniano.
        \item \( G \) es hamiltoniano, pero no euleriano.
    \end{enumerate}
\end{ejercicio}

\vspace{0.5cm}

\begin{definition}
    Dados \( m, n \geq 1 \), el \textbf{grafo bipartito completo} \( K_{m,n} \) es el grafo bipartito con bipartición \( A, B \) tal que \( |A| = m \), \( |B| = n \), y cada vértice de \( A \) es adyacente a cada vértice de \( B \) (en particular, \( |E(K_{m,n})| = mn \)).
\end{definition}

\begin{ejercicio}{3}
    Para \( n > 1 \), probar que \( K_{n,n} \) posee \( \dfrac{(n-1)!\,n!}{2} \) ciclos hamiltonianos.
\end{ejercicio}

\vspace{0.5cm}

\begin{definition}
    Un \textbf{camino hamiltoniano} de un grafo \( G \) es un camino \( P \) que contiene a todos los vértices de \( G \).
\end{definition}

\begin{ejercicio}{4}
    Todo grafo con un ciclo hamiltoniano posee un camino hamiltoniano, pero no vale la recíproca (basta tomar un grafo camino como contraejemplo). Probar que un grafo \( G \) posee un camino hamiltoniano si y sólo si el grafo \( G' \) obtenido de \( G \) agregándole un vértice nuevo \( u \) y haciéndolo adyacente a todos los demás vértices posee un ciclo hamiltoniano.
\end{ejercicio}

\vspace{0.5cm}

\begin{ejercicio}{5}
    Un ratón intenta comerse un cubo de 3x3x3 de queso. Empieza en una esquina y se come cada vez un cubo de 1x1x1 antes de pasar a uno adyacente (decimos que dos cubos son adyacentes cuando comparten cara). ¿Puede comerse el cubo central en último lugar?
\end{ejercicio}

\vspace{0.5cm}

\begin{ejercicio}{6}
    Elena y Dora invitan a 10 amigos a cenar. En este grupo de doce, todos conocen a al menos seis personas. Demostrar que pueden sentarse los doce alrededor de una mesa circular de modo que todo comensal conozca a las dos personas que están sentadas a su lado.
\end{ejercicio}

\vspace{0.5cm}

\begin{ejercicio}{7}
    Probar que, para todo \( n > 2 \), existe un grafo simple conexo no hamiltoniano tal que todo vértice tiene grado al menos \( \left\lfloor \dfrac{n-1}{2} \right\rfloor \). Vemos así que el grado mínimo se puede acercar mucho a \( \dfrac{n}{2} \) sin que el grafo sea hamiltoniano.
\end{ejercicio}

\vspace{0.5cm}

\begin{ejercicio}{8}
    Sea \( n \) un entero mayor o igual que 3. Probar que el máximo número posible de aristas de un grafo simple no hamiltoniano con \( n \) vértices es \( \dbinom{n-1}{2} + 1 \).

    (Sugerencia: probar que, en un grafo con más aristas, \( d(u) + d(v) \geq n \) para todo par de vértices \( u \) y \( v \) no adyacentes. Alternativamente, probar por inducción que todo grafo con \( \dbinom{n-1}{2} + 2 \) aristas es hamiltoniano. Para esto último ayuda hallar cuáles son los posibles valores del máximo grado de un vértice).
\end{ejercicio}

\end{document}

\clearpage

\section{Práctica VI}
\begin{ejercicio}{1}
	Determinar si es verdadero o falso, y justificar: si \( u \) y \( v \) son los únicos vértices de grado impar en un grafo, entonces \( G \) contiene un \( uv \)-camino.
\end{ejercicio}
\begin{proof}
	Notemos primero que \( u \) y \( v \) deben pertenecer a la misma componente conexa de \( G \): en cualquier componente conexa, la suma de los grados de sus vértices es par, por lo que la cantidad de vértices de grado impar dentro de ella también es par. Como \( u \) y \( v \) son los \emph{únicos} vértices de grado impar de \( G \), una componente que contenga a uno de ellos no puede tener una cantidad impar de vértices de grado impar quedando el otro afuera; luego ambos están en la misma componente \( H \) (las demás componentes, si las hay, tienen todos sus vértices de grado par).

	Trabajemos entonces dentro de \( H \). Agreguemos la arista \( uv \) a \( H \) y llamemos \( H' \) al grafo resultante. En \( H' \), todos los vértices tienen grado par y \( H' \) sigue siendo conexo, por lo que \( H' \) tiene un circuito euleriano. Este circuito euleriano debe pasar por la arista \( uv \), pero podemos mirar la parte del circuito que va de \( u \) a \( v \) sin pasar por \( uv \) o, en el otro caso, el que va de \( v \) a \( u \) sin pasar por esa arista, obteniendo así un \( uv \)-camino (que también es un \( uv \)-camino en \( G \), ya que \( H \) es subgrafo de \( G \)).
\end{proof}

\vspace{0.5cm}

\begin{ejercicioEntregar}{2}
	Probar si las siguientes son secuencias de grados de un grafo simple. En caso afirmativo, construir un grafo con esos grados. Si no, dar una prueba de imposibilidad.
	\begin{itemize}
		\item \( 5, 5, 4, 3, 2, 2, 2, 1 \)
		\item \( 5, 5, 4, 4, 2, 2, 1, 1 \)
		\item \( 5, 5, 5, 3, 2, 2, 1, 1 \)
		\item \( 5, 5, 5, 4, 2, 1, 1, 1 \)
	\end{itemize}
\end{ejercicioEntregar}
\begin{proof}
	Recordemos el \textbf{algoritmo} visto en la teoría para determinar si una secuencia de enteros no negativos \( d_1, \ldots, d_n \) es la secuencia de grados de un grafo simple. Sea \( d_1 \geq d_2 \geq \cdots \geq d_n \). Si \( d_1 = 0 \), entonces todos los \( d_i = 0 \) y la secuencia es gráfica. Si \( d_1 > 0 \), eliminamos el primer elemento \( d_1 \) y le restamos 1 a los siguientes \( d_1 \) elementos. Luego reordenamos la secuencia resultante y repetimos el proceso hasta que lleguemos a una secuencia de ceros o a una secuencia que no sea válida (por ejemplo, si algún elemento se vuelve negativo).

	Hagamos esto para cada una de las secuencias dadas:
	\begin{itemize}
		\item \textbf{Secuencia 1:} \( 5, 5, 4, 3, 2, 2, 2, 1 \)
		      \[
			      5, 5, 4, 3, 2, 2, 2, 1 \to 4, 3, 2, 1, 1, 2, 1 \to 4, 3, 2, 2, 1, 1, 1 \to 2, 1, 1, 0, 1, 1 \to 2, 1, 1, 1, 1, 0
		      \]
		      \[
			      \to 0, 0, 1, 1, 0 \to 1, 1, 0, 0, 0 \to 0, 0, 0, 0 \implies \text{Secuencia válida.}
		      \]

		\item \textbf{Secuencia 2:} \( 5, 5, 4, 4, 2, 2, 1, 1 \)
		      \[
			      5, 5, 4, 4, 2, 2, 1, 1 \to 4, 3, 3, 1, 1, 1, 1 \to 2, 2, 0, 0, 1, 1 \to 2, 2, 1, 1, 0, 0
		      \]
		      \[
			      \to 1, 0, 1, 0, 0 \to 1, 1, 0, 0, 0 \to 0, 0, 0, 0 \implies \text{Secuencia válida.}
		      \]

		\item \textbf{Secuencia 3:} \( 5, 5, 5, 3, 2, 2, 1, 1 \)
		      \[
			      5, 5, 5, 3, 2, 2, 1, 1 \to 4, 4, 2, 1, 1, 1, 1 \to 3, 1, 0, 0, 1, 1 \to 3, 1, 1, 1, 0, 0
		      \]
		      \[
			      \to 0, 0, 0, 0, 0 \implies \text{Secuencia válida.}
		      \]

		\item \textbf{Secuencia 4:} \( 5, 5, 5, 4, 2, 1, 1, 1 \)
		      \[
			      5, 5, 5, 4, 2, 1, 1, 1 \to 4, 4, 3, 1, 0, 1, 1 \to 4, 4, 3, 1, 1, 1, 0 \to 3, 2, 0, 0, 1, 0
		      \]
		      \[
			      \to 3, 2, 1, 0, 0, 0 \to 1, 0, -1, 0, 0 \implies \text{Secuencia \textbf{NO} válida, aparece un grado negativo}
		      \]
	\end{itemize}

	\vspace{0.5cm}
	A continuación, construimos los grafos simples correspondientes para las secuencias válidas. Nombramos los vértices \(v_1, \ldots, v_8\) donde el subíndice respeta el orden original de la secuencia de mayor a menor grado.

	\begin{center}
		\begin{tikzpicture}[
				vertex/.style={circle, draw, minimum size=8mm, inner sep=0pt, font=\small},
				thick]

			\node at (0, 3) {\textbf{Grafo 1}};
			% Grafo 1: 5, 5, 4, 3, 2, 2, 2, 1
			\begin{scope}[shift={(0,0)}]
				\foreach \i/\ang in {1/90, 2/45, 3/0, 4/-45, 5/-90, 6/-135, 7/-180, 8/135} {
						\node[vertex] (g1v\i) at (\ang:2) {$v_{\i}$};
					}
				% Aristas Grafo 1
				\draw (g1v1) -- (g1v2); \draw (g1v1) -- (g1v3); \draw (g1v1) -- (g1v4); \draw (g1v1) -- (g1v5); \draw (g1v1) -- (g1v6);
				\draw (g1v2) -- (g1v3); \draw (g1v2) -- (g1v4); \draw (g1v2) -- (g1v5); \draw (g1v2) -- (g1v6);
				\draw (g1v3) -- (g1v7); \draw (g1v3) -- (g1v8);
				\draw (g1v4) -- (g1v7);
			\end{scope}

			\node at (6, 3) {\textbf{Grafo 2}};
			% Grafo 2: 5, 5, 4, 4, 2, 2, 1, 1
			\begin{scope}[shift={(6,0)}]
				\foreach \i/\ang in {1/90, 2/45, 3/0, 4/-45, 5/-90, 6/-135, 7/-180, 8/135} {
						\node[vertex] (g2v\i) at (\ang:2) {$v_{\i}$};
					}
				% Aristas Grafo 2
				\draw (g2v1) -- (g2v2); \draw (g2v1) -- (g2v3); \draw (g2v1) -- (g2v4); \draw (g2v1) -- (g2v5); \draw (g2v1) -- (g2v6);
				\draw (g2v2) -- (g2v3); \draw (g2v2) -- (g2v4); \draw (g2v2) -- (g2v5); \draw (g2v2) -- (g2v6);
				\draw (g2v3) -- (g2v4); \draw (g2v3) -- (g2v7);
				\draw (g2v4) -- (g2v8);
			\end{scope}

			\node at (12, 3) {\textbf{Grafo 3}};
			% Grafo 3: 5, 5, 5, 3, 2, 2, 1, 1
			\begin{scope}[shift={(12,0)}]
				\foreach \i/\ang in {1/90, 2/45, 3/0, 4/-45, 5/-90, 6/-135, 7/-180, 8/135} {
						\node[vertex] (g3v\i) at (\ang:2) {$v_{\i}$};
					}
				% Aristas Grafo 3
				\draw (g3v1) -- (g3v2); \draw (g3v1) -- (g3v3); \draw (g3v1) -- (g3v4); \draw (g3v1) -- (g3v5); \draw (g3v1) -- (g3v6);
				\draw (g3v2) -- (g3v3); \draw (g3v2) -- (g3v4); \draw (g3v2) -- (g3v5); \draw (g3v2) -- (g3v7);
				\draw (g3v3) -- (g3v4); \draw (g3v3) -- (g3v6); \draw (g3v3) -- (g3v8);
			\end{scope}
		\end{tikzpicture}
	\end{center}
\end{proof}

\vspace{0.5cm}

\begin{ejercicio}{3}
	Se tiene una liga con dos divisiones de 13 equipos cada una. Determinar si es posible programar la temporada de manera que cada equipo juegue nueve partidos con equipos de su división y cuatro partidos con equipos de la otra división.
\end{ejercicio}
\begin{proof}
	Si separamos cada división y vemos los partidos que se juegan dentro de cada división, tenemos que cada equipo juega 9 partidos con equipos de su división. Esto quiere decir que, si pensamos cada equipo como un vértice, tenemos que armar el grafo con la siguiente secuencia de números: \[
		9, 9, 9, 9, 9, 9, 9, 9, 9, 9, 9, 9, 9
	\]
	Cuya suma total es \( 9 \cdot 13 = 117 \), que es impar. Por lo tanto, no es posible que los equipos jueguen 9 partidos con equipos de su división y, por lo tanto, no es posible programar la temporada de la manera que se pide.
\end{proof}

\vspace{0.5cm}

\begin{ejercicioEntregar}{4}
	Probar que todo grafo simple con al menos dos vértices tiene dos vértices de igual grado ¿Vale lo mismo para grafos sin loops con aristas múltiples?
\end{ejercicioEntregar}
\begin{proof}
	Sea \( G \) un grafo simple con \( n \geq 2 \) vértices. Los posibles grados de los vértices en un grafo simple con \( n \) vértices son \( 0, 1, 2, \ldots, n-1 \). Esto nos da \( n \) posibles grados. Supongamos que los \( n \) vértices tienen todos grados distintos. Si aplicamos el algoritmo, nos queda:\[
		n-1, n-2, n-3, \ldots, 2, 1, 0 \to n-3, n-4, \ldots, 1, -1
	\]
	Luego, no es posible armar un grafo de estas características.

	No vale lo mismo para grafos sin loops con aristas múltiples. Por ejemplo, el grafo de 3 vértices con una simple y una arista múltiple entre dos de ellos tiene grados \( 1, 2, 3 \) y no hay dos vértices con el mismo grado.
\end{proof}

\vspace{0.5cm}

\begin{ejercicio}{5}
	Sean \( d_1, \ldots, d_n \) enteros tales que \( d_1 \geq \cdots \geq d_n \geq 0 \). Probar que existe un grafo sin loops (se permiten aristas múltiples) con secuencia de grados \( d_1, \ldots, d_n \) si y sólo si la suma de estos números es par y \( d_1 \leq d_2 + \cdots + d_n \).
\end{ejercicio}
\begin{proof}
	Sean \( d_1, \ldots, d_n \) enteros tales que \( d_1 \geq \cdots \geq d_n \geq 0 \).
	\begin{itemize}
		\item[\(\Longrightarrow\))] Supongamos que existe un grafo \( G \) sin loops con secuencia de grados \( d_1, \ldots, d_n \). La suma de los grados siempre es par. Luego,
		      \[
			      \sum_{i=1}^n d_i = 2 |E(G)|
		      \]
		      Como no hay loops, el grado de cualquier vértice a lo sumo puede ser igual a la cantidad total de aristas en el grafo, por lo que \( d_1 \le |E(G)| \). Sustituyendo esto, tenemos:
		      \[
			      d_1 \le \sum_{i=1}^n \frac{d_i}{2} \implies 2d_1 \le \sum_{i=1}^n d_i \implies 2d_1 \le d_1 + \sum_{i=2}^n d_i \implies d_1 \le d_2 + \cdots + d_n
		      \]

		\item[\(\Longleftarrow\))] Sea \( d_1 \le d_2 + \cdots + d_n \) y supongamos que \( \sum_{i=1}^n d_i \) es par.
		      Llamemos \( S = \sum_{i=1}^n d_i \). Como \( S \) es par, podemos escribir \( S = 2m \).
		      Por hipótesis, \( d_1 \le \sum_{i=2}^n d_i \). Sumando \( d_1 \) a ambos lados obtenemos \( 2d_1 \le S = 2m \), lo que implica que \( d_1 \le m \).

		      Vamos a construir un grafo \( G \). Consideremos un conjunto de \( 2m \) puntos, numerados del \( 1 \) al \( 2m \).
		      Asignamos estos puntos a los vértices de manera secuencial en bloques contiguos:
		      \begin{itemize}
			      \item Al vértice \( v_1 \) le asignamos los primeros \( d_1 \) puntos.
			      \item Al vértice \( v_2 \) le asignamos los siguientes \( d_2 \) puntos.
			      \item Y así sucesivamente; en general, a cada vértice \( v_i \) le asignamos un bloque de \( d_i \) puntos.
		      \end{itemize}

		      Para formar las aristas de \( G \), unimos el punto \( j \) con el punto \( j+m \) para todo \( 1 \le j \le m \).

		      Es evidente que cada vértice \( v_i \) tendrá grado exactamente \( d_i \), puesto que le asignamos \( d_i \) puntos y cada punto forma exactamente un extremo de una arista.

		      Resta verificar que el grafo resultante no tiene loops. Un loop se formaría si conectamos el punto \( j \) con el punto \( j+m \) y ambos pertenecen al mismo vértice \( v_i \). Sin embargo, la cantidad de puntos asignados al vértice \( v_i \) es \( d_i \). Como la secuencia está ordenada de forma decreciente, el vértice con más puntos es \( v_1 \), el cual tiene \( d_1 \) puntos.
		      Dado que \( d_1 \le m \), la distancia máxima entre los números de dos puntos que pertenecen al mismo vértice es estrictamente menor que \( m \).
		      Por lo tanto, el punto \( j \) y el punto \( j+m \) nunca pueden pertenecer al mismo vértice. Esto garantiza que no se forman loops, concluyendo la demostración.
	\end{itemize}
\end{proof}

\vspace{0.5cm}

\begin{ejercicio}{6}
	Determinar si es verdadero o falso y justificar: todo digrafo simple (permitiendo loops) con \( n \) vértices (\( n > 1 \)) tiene dos vértices distintos con el mismo grado de salida o dos vértices distintos con el mismo grado de llegada ¿Cambia la respuesta si pensamos a los digrafos simples en el sentido más restrictivo de que los loops no están permitidos?
\end{ejercicio}
\begin{proof}
	\emph{Claim}: La afirmación es \textbf{falsa} en ambos casos.

	\begin{center}
		\begin{tikzpicture}[
				vertex/.style={circle, draw, minimum size=8mm, inner sep=0pt, font=\small}, ->, >=stealth, thick]

			\node[vertex] (v1) at (0, 0) {$v_1$};
			\node[vertex] (v2) at (4, 0) {$v_2$};
			\node[vertex] (v3) at (2, 3) {$v_3$};

			\draw (v3) -- (v2);
			\draw (v3) -- (v1);
			\draw (v2) -- (v1);

			% Etiquetas de grados
			\node[align=left, font=\footnotesize, right=0.3cm] at (v2) {$d^+(v_2)=1$\\$d^-(v_2)=1$};
			\node[align=right, font=\footnotesize, left=0.3cm] at (v1) {$d^+(v_1)=0$\\$d^-(v_1)=2$};
			\node[align=center, font=\footnotesize, above=0.3cm] at (v3) {$d^+(v_3)=2$\\$d^-(v_3)=0$};

		\end{tikzpicture}
	\end{center}

	Los grados de salida son \( d^+(v_1)=0 \), \( d^+(v_2)=1 \), \( d^+(v_3)=2 \) (todos distintos).
	Los grados de llegada son \( d^-(v_1)=2 \), \( d^-(v_2)=1 \), \( d^-(v_3)=0 \) (todos distintos).

	Como el grafo no utiliza loops, los grados de salida y llegada siguen siendo distintos para cada vértice. Por lo tanto, la afirmación es falsa en ambos casos.
\end{proof}

\vspace{0.5cm}

\begin{definition}
	Un \textbf{camino dirigido} en un digrafo \( D \) es una secuencia de vértices \( v_1, v_2, \ldots, v_k \) tal que para cada \( i = 1, 2, \ldots, k-1 \), existe una arista dirigida de \( v_i \) a \( v_{i+1} \).
\end{definition}

\begin{definition}
	Un digrafo se dice \textbf{fuertemente conexo} si para cada par de vértices \( u \) y \( v \), existe un camino dirigido de \( u \) a \( v \) y otro de \( v \) a \( u \).
\end{definition}

\begin{ejercicioEntregar}{7}
	Dos vértices de un digrafo \( x \) e \( y \) se dicen conectados si existe un camino dirigido de \( x \) a \( y \) y otro de \( y \) a \( x \). Probar que la relación de conexión en un digrafo \( D \) es de equivalencia y que sus clases son las componentes fuertemente conexas de \( D \).
\end{ejercicioEntregar}
\begin{proof}
	Para ver que es una relación de equivalencia, debemos verificar que es reflexiva, simétrica y transitiva. Llamemos \( \sim \) a esta relación de conexión, es decir, \( x \sim y \) si y sólo si existe un camino dirigido de \( x \) a \( y \) y otro de \( y \) a \( x \).

	\begin{itemize}
		\item \textbf{Reflexiva (\( x \sim x \)):} Para cualquier vértice \( x \in V(D) \), la secuencia formada por el único vértice \( x \) constituye un camino dirigido de longitud 0 (camino trivial o vacío) de \( x \) a \( x \). Por lo tanto, \( x \sim x \).

		\item \textbf{Simétrica (\( x \sim y \implies y \sim x \)):} Si \( x \sim y \), entonces por definición existe un camino dirigido de \( x \) a \( y \) y otro de \( y \) a \( x \). Esto es equivalente a decir que existe un camino de \( y \) a \( x \) y otro de \( x \) a \( y \), por lo que \( y \sim x \).

		\item \textbf{Transitiva (\( x \sim y \) e \( y \sim z \implies x \sim z \)):} Si \( x \sim y \), existen caminos dirigidos \( P_{xy} \) (de \( x \) a \( y \)) y \( P_{yx} \) (de \( y \) a \( x \)). Como \( y \sim z \), existen caminos dirigidos \( P_{yz} \) y \( P_{zy} \). Concatenando \( P_{xy} \) con \( P_{yz} \) obtenemos un recorrido dirigido de \( x \) a \( z \), del cual siempre podemos extraer un camino dirigido de \( x \) a \( z \). Análogamente, concatenando \( P_{zy} \) con \( P_{yx} \) obtenemos un recorrido del cual extraemos un camino dirigido de \( z \) a \( x \). Por lo tanto, \( x \sim z \).
	\end{itemize}

	Luego, la relación de conexión \( \sim \) es de equivalencia.

	Veamos ahora que las clases de equivalencia bajo esta relación son precisamente las componentes fuertemente conexas del digrafo \( D \).

	Por un lado, la clase de equivalencia de un vértice \( x \), denotada \( [x] \), es el conjunto de todos los vértices \( y \) tales que \( x \sim y \). Por definición de \( \sim \), todos los vértices en \( [x] \) son mutuamente alcanzables, formando un conjunto fuertemente conexo.

	Además, las clases de equivalencia son subconjuntos \textbf{maximales} respecto a esta propiedad: si existiera un vértice \( w \notin [x] \) que fuera mutuamente alcanzable con los vértices de \( [x] \), en particular tendríamos \( x \sim w \), lo que implicaría \( w \in [x] \), llegando a un absurdo.

	Dado que una componente fuertemente conexa se define justamente como un subgrafo fuertemente conexo maximal, concluimos que las clases de equivalencia de \( \sim \) coinciden exactamente con las componentes fuertemente conexas de \( D \).
\end{proof}

\vspace{0.5cm}

\begin{ejercicioEntregar}{8}
	Probar que en todo digrafo existe una componente fuertemente conexa sin aristas que ingresen a ella y que existe una componente fuertemente conexa sin aristas que salgan de ella.
\end{ejercicioEntregar}
\begin{proof}
	Sea \( D \) un digrafo (asumimos finito) y consideremos su grafo de componentes fuertemente conexas, \( D^* \). En \( D^* \), cada vértice representa una componente fuertemente conexa de \( D \), y hay una arista dirigida de la componente \( C_i \) a la componente \( C_j \) si existe al menos una arista en \( D \) desde algún vértice en \( C_i \) hacia algún vértice en \( C_j \).

	Si el grafo \( D^* \) tuviera un ciclo dirigido, esto implicaría que todos los vértices de las componentes involucradas en el ciclo serían mutuamente alcanzables en \( D \). Pero entonces todas ellas formarían una única gran componente fuertemente conexa, lo que contradice la definición de las componentes originales. Por lo tanto, \( D^* \) es un digrafo acíclico.

	Como \( D \) es finito, \( D^* \) también lo es. Veamos que en todo digrafo acíclico finito existe al menos un vértice sin aristas entrantes y al menos uno sin aristas salientes.

	Consideremos un camino dirigido de longitud máxima en \( D^* \), digamos \( P = C_1, C_2, \ldots, C_k \).
	\begin{itemize}
		\item El vértice inicial \( C_1 \) no puede tener aristas que ingresen a él. Si recibiera una arista de algún \( C_i \) perteneciente a \( P \), se formaría un ciclo (absurdo). Si recibiera una arista de un vértice externo a \( P \), el camino podría extenderse hacia atrás, contradiciendo que \( P \) es de longitud máxima. Por lo tanto, \( C_1 \) es la componente buscada sin aristas entrantes.
		\item Por un razonamiento análogo, el vértice final \( C_k \) no puede tener aristas que salgan de él (de lo contrario se formaría un ciclo o el camino se extendería hacia adelante). Luego, \( C_k \) es la componente buscada sin aristas salientes.
	\end{itemize}

	Esto demuestra la existencia de ambas componentes en el digrafo original \( D \).
\end{proof}

\vspace{0.5cm}

\begin{ejercicio}{9}
	Sea \( G \) un digrafo simple con \( n \) vértices y sin ciclos dirigidos. Probar que existe un orden \( v_1, \ldots, v_n \) de los vértices de \( G \) de modo que, para todo par de valores distintos \( i, j \) entre \( 1 \) y \( n \), si \( v_i v_j \in E(G) \) entonces \( i \leq j \).
\end{ejercicio}
\begin{proof}
	Procedemos por inducción sobre \( n \), la cantidad de vértices de \( G \).

	\paragraph{Caso base:} Si \( n = 1 \), \( G \) tiene un único vértice \( v_1 \) y ninguna arista. El orden \( v_1 \) satisface la condición trivialmente.

	\paragraph{Paso inductivo:} Supongamos que la afirmación es cierta para todo digrafo simple sin ciclos dirigidos de \( n-1 \) vértices (\( n \ge 2 \)).

	Sea \( G \) un digrafo simple sin ciclos dirigidos de \( n \) vértices. Sea \( S \) el conjunto de los vértices de \( G \) con grado de entrada igual a cero.

	Primero, demostremos que \( S \neq \varnothing \). Supongamos por el absurdo que \( S = \varnothing \implies \) todo vértice de \( G \) tiene al menos una arista entrante. Si partimos de un vértice cualquiera y retrocedemos por las aristas entrantes, podemos construir un camino infinito. Dado que \( G \) tiene una cantidad finita de vértices, en algún momento deberemos visitar un vértice por segunda vez. Esto forma un ciclo dirigido, pero \( G \) es acíclico \( \implies S \neq \varnothing \).

	Tomemos \( v \in S \) y lo fijamos como el primer elemento de la secuencia, \( v_1 = v \). Consideremos el subgrafo inducido \( G' = G - \{v_1\} \). Es claro que \( G' \) es un digrafo simple sin ciclos dirigidos y tiene \( n-1 \) vértices.

	Por \textbf{H.I}, existe un orden \( v_2, \ldots, v_n \) de los vértices de \( G' \) tal que para todo par de índices \( i, j \in \{2, \ldots, n\} \), si \( v_i v_j \in E(G') \), entonces \( i \le j \).

	Para finalizar, comprobemos que el orden \( v_1, v_2, \ldots, v_n \) es válido para la totalidad de \( G \). Tomemos una arista cualquiera \( v_i v_j \in E(G) \) y evaluemos los casos posibles:

	\begin{itemize}
		\item \textbf{Si \( i, j \ge 2 \):} estamos en la \textbf{H.I} y vale.
		\item \textbf{Si \( i = 1 \):} La arista es de la forma \( v_1 v_j \). Como \( j \in \{2, \ldots, n\} \implies 1 < j \implies i \le j \).
		\item \textbf{Si \( j = 1 \):} No existen aristas de la forma \( v_i v_1 \) pues \( v_1 \in S \).
	\end{itemize}

	En todos los casos posibles, \( v_i v_j \in E(G) \implies i \le j \). Esto completa la demostración.
\end{proof}

\vspace{0.5cm}

\begin{definition}
	Un \textbf{torneo} es un digrafo cuyo grafo subyacente es un grafo completo. En otras palabras, el digrafo obtenido de un grafo completo al darle orientación a sus aristas.
\end{definition}

\begin{ejercicio}{10}
	Probar que existe un torneo de \( n \) vértices con grado de salida igual al grado de llegada en cada vértice si y sólo si \( n \) es impar.
\end{ejercicio}
\begin{proof}
	\(\Longrightarrow\)) Sea \( G \) un torneo de \( n \) vértices tal que en cada vértice el grado de salida es igual al grado de llegada. En cualquier torneo, \(\forall v \in V(G) \) vale que \( d^+(v) + d^-(v) = n-1 \). Como por hipótesis \( d^+(v) = d^-(v) \), obtenemos \( 2d^+(v) = n-1 \implies n \) es impar.

	\(\Longleftarrow)\) Supongamos que \( n \) es impar. Queremos construir un torneo de \( n \) vértices tal que \( d^+(v) = d^-(v) \quad \forall v \in V(G) \). Sea \( k = \frac{n-1}{2} \in \N \).

	Sea el conjunto de vértices \( V = \{0, 1, \ldots, n-1\} \). Trazamos una arista dirigida de \( i \) a \( j \iff (j - i) \pmod n \in \{1, 2, \ldots, k\} \).

	Veamos que esto define un torneo. Dados dos vértices distintos \( i, j \in V \), como \( i \not\equiv j \pmod n \), ambos restos \( (j-i) \bmod n \) y \( (i-j) \bmod n \) son no nulos y pertenecen a \( \{1, \ldots, n-1\} \); además, al ser uno el opuesto del otro módulo \( n \), su suma es exactamente \( n \). Si ambos fueran \( \le k \), la suma sería a lo sumo \( 2k < n \); si ambos fueran \( \ge k+1 \), la suma sería al menos \( 2k+2 > n \). Luego, exactamente uno de los dos restos es \( \le k \), es decir, se traza exactamente una arista dirigida entre cada par de vértices, cumpliendo la definición de torneo.

	Finalmente, por construcción, para cada vértice \( i \) hay exactamente \( k \) valores distintos de \( j \) hacia los que apunta. Por lo tanto, el grado de salida de todo vértice es \( k = \frac{n-1}{2} \). Como cada vértice está conectado a \( n-1 \) vértices en total, su grado de entrada también es \( n - 1 - k = k \), cumpliéndose así lo requerido.
\end{proof}

\vspace{0.5cm}

\begin{definition}
	Un \textbf{camino dirigido hamiltoniano} en un digrafo \( D \) es un camino dirigido que visita cada vértice de \( D \) exactamente una vez.
\end{definition}

\begin{ejercicioEntregar}{11}
	Probar que todo torneo posee un camino dirigido hamiltoniano.
\end{ejercicioEntregar}
\begin{proof}
	Sea \( G \) un torneo con \( n \) vértices. Procedemos por inducción sobre \( n \).

	\paragraph{Caso base:} Si \( n = 1 \), el único vértice forma un camino hamiltoniano trivial.

	\paragraph{Paso inductivo:} Supongamos que vale para todo torneo con \( n-1 \) vértices. Sea \( G \) un torneo con \( n \) vértices. Tomemos un vértice \( v \in V(G) \) y consideremos el subgrafo inducido \( G' = G - \{v\} \), que es un torneo de \( n-1 \) vértices. Por \textbf{H.I}, existe un camino dirigido hamiltoniano en \( G' \), digamos \( P = v_1, v_2, \ldots, v_{n-1} \).

	Para extender este camino a todo \( G \), debemos insertar el vértice \( v \). Como \( G \) es un torneo, entre \( v \) y cualquier \( v_k \) existe exactamente una arista dirigida. Analizamos tres casos posibles para la inserción:

	\begin{itemize}
		\item Si \( v v_1 \in E(G) \), entonces el camino \( v, v_1, v_2, \ldots, v_{n-1} \) es un camino dirigido hamiltoniano en \( G \).

		\item Si \( v_{n-1} v \in E(G) \), entonces el camino \( v_1, v_2, \ldots, v_{n-1}, v \) es un camino dirigido hamiltoniano en \( G \).

		\item Veamos qué ocurre si \( v v_1 \notin E(G) \) y \( v_{n-1} v \notin E(G) \). Sea \( i \) el \textit{máximo} índice en el conjunto \( \{1, \ldots, n-1\} \) tal que \( v_i v \in E(G) \).
		      Sabemos que \( i \) existe y es al menos \( 1 \) (pues \( v_1 v \in E(G) \)), y sabemos que \( i < n-1 \) (pues \( v_{n-1} v \notin E(G) \)).

		      Consideremos el vértice siguiente en el camino, \( v_{i+1} \). Como \( i \) es el máximo índice que apunta a \( v \), la arista entre \( v_{i+1} \) y \( v \) no puede estar dirigida hacia \( v \). Por ser \( G \) un torneo, necesariamente sale de \( v \), es decir, \( v v_{i+1} \in E(G) \).

		      De esta manera podemos insertar \( v \) justo entre ellos, formando el camino dirigido hamiltoniano \( v_1, \ldots, v_i, v, v_{i+1}, \ldots, v_{n-1} \).
	\end{itemize}

	En todos los casos es posible construir el camino hamiltoniano, lo que completa la demostración.
\end{proof}

\vspace{0.5cm}

\begin{definition}
	Un \textbf{ciclo dirigido hamiltoniano} en un digrafo \( D \) es un ciclo dirigido que visita cada vértice de \( D \) exactamente una vez.
\end{definition}

\begin{ejercicio}{12}
	Probar que un torneo posee un ciclo hamiltoniano si y sólo si es fuertemente conexo.
\end{ejercicio}
\begin{proof}
	(\(\Longrightarrow\)) Sea \( G \) un torneo con un ciclo hamiltoniano. Sea \( C = v_1, v_2, \ldots, v_n, v_1 \) dicho ciclo. Para cualquier par de vértices \( v_i, v_j \), podemos ir de \( v_i \) a \( v_j \) siguiendo el ciclo en la dirección adecuada. Por lo tanto, existe un camino dirigido de \( v_i \) a \( v_j \) y otro de \( v_j \) a \( v_i \), lo que implica que \( G \) es fuertemente conexo.

	(\(\Longleftarrow\)) Sea \( G \) un torneo fuertemente conexo. Demostraremos que contiene un ciclo hamiltoniano tomando un ciclo de longitud máxima y probando por el absurdo que debe incluir a todos los vértices.

	Como \( G \) es fuertemente conexo (y \( n \ge 3 \), ya que para \( n=1 \) es trivial y \( n=2 \) no puede ser fuertemente conexo), cada vértice tiene grado de entrada y salida mayor a cero, lo que garantiza la existencia de al menos un ciclo.
	Sea \( C = v_1, v_2, \ldots, v_k, v_1 \) un ciclo dirigido de longitud máxima \( k \). Supongamos por el absurdo que \( k < n \).

	Sea \( W = V(G) \setminus V( C ) \) el conjunto de vértices que no están en el ciclo. Evaluamos las conexiones entre \( W \) y \( C \). Existen dos posibilidades para un vértice \( u \in W \):

	\paragraph{Caso 1:} \( u \) tiene al menos una arista hacia \( C \) y al menos una arista recibida desde \( C \).
	Por un razonamiento anáogo al ejercicio anterior, al recorrer los vértices de \( C \) en orden, la dirección de las aristas entre \( C \) y \( u \) debe cambiar en algún momento de salir de \( C \) a entrar a \( C \). Es decir, existe un índice \( i \) tal que \( v_i u \in E(G) \) y \( u v_{i+1} \in E(G) \).
	Si insertamos \( u \), formamos el ciclo \( v_1, \ldots, v_i, u, v_{i+1}, \ldots, v_k, v_1 \) de longitud \( k+1 \), lo que contradice que \( C \) era de longitud máxima.

	\paragraph{Caso 2:} Ningún vértice de \( W \) cumple el Caso 1.
	Entonces, \( W \) se divide estrictamente en dos subconjuntos: \( W_{in} \), formado por los vértices que \textit{sólo} reciben aristas desde \( C \) (y por lo tanto no envían ninguna arista hacia \( C \)); y \( W_{out} \), formado por los vértices que \textit{sólo} envían aristas hacia \( C \) (y por lo tanto no reciben ninguna arista desde \( C \)).

	Veamos primero que ambos subconjuntos son no vacíos. Si \( W_{out} = \varnothing \), entonces todo vértice de \( W \) pertenece a \( W_{in} \), y por definición ninguno de ellos envía una arista hacia \( C \); como tampoco hay otra forma de salir de \( W \) (las únicas aristas entre \( C \) y \( W \) son las que van de \( C \) a \( W_{in} \)), ningún vértice de \( W \) podría alcanzar jamás a un vértice de \( C \), contradiciendo que \( G \) es fuertemente conexo (recordar que \( W \neq \varnothing \) pues \( k < n \)). Por lo tanto \( W_{out} \neq \varnothing \). Simétricamente, si \( W_{in} = \varnothing \), entonces todo vértice de \( W \) pertenece a \( W_{out} \) y ninguno recibe una arista desde \( C \), por lo que ningún vértice de \( C \) podría alcanzar a un vértice de \( W \), contradiciendo nuevamente que \( G \) es fuertemente conexo. Luego \( W_{in} \neq \varnothing \).

	Afirmamos ahora que existe al menos un par \( y \in W_{in} \), \( x \in W_{out} \) con \( yx \in E(G) \). En efecto, supongamos por el absurdo que para todo \( y \in W_{in} \) y todo \( x \in W_{out} \) la arista va en el sentido \( xy \) (esto cubre todos los casos, pues \( G \) es un torneo y entre \( y \) y \( x \) hay exactamente una arista dirigida). Entonces ningún vértice de \( W_{in} \) tiene aristas salientes hacia \( C \) (por definición de \( W_{in} \)) ni hacia \( W_{out} \) (por la suposición), de modo que sus únicas aristas salientes posibles quedan dentro del propio \( W_{in} \). Esto implica que ningún vértice de \( W_{in} \) puede alcanzar jamás a un vértice de \( C \), contradiciendo que \( G \) es fuertemente conexo. Luego existen \( y \in W_{in} \) y \( x \in W_{out} \) con \( yx \in E(G) \).

	Pero en este escenario, sabiendo que \( v_1 \) apunta a \( y \) (pues \( y \in W_{in} \)) y que \( x \) apunta a \( v_2 \) (pues \( x \in W_{out} \)), podemos formar el ciclo \( v_1, y, x, v_2, v_3, \ldots, v_k, v_1 \). Este nuevo ciclo tiene longitud \( k+2 \), lo que vuelve a contradecir la maximidad de \( C \).

	En ambos casos llegamos a un absurdo producto de suponer que \( k < n \). Por lo tanto, \( k = n \) y el ciclo máximo \( C \) es un ciclo hamiltoniano.
\end{proof}

\clearpage

\end{document}
